% Luc Destombe --- licence CC BY-NC-SA 4.0
% https://creativecommons.org/licenses/by-nc-sa/4.0/deed.fr
% Fichier autonome : compiler avec pdflatex, deux fois (signets, nombre de pages).
\documentclass[11pt,a4paper]{article}
\makeatletter
\newif\iflycee@unecolonne
\lycee@unecolonnetrue
\newif\iflycee@palatino
\lycee@palatinotrue

\RequirePackage[utf8]{inputenc}
\RequirePackage[T1]{fontenc}
\RequirePackage[french]{babel}
\RequirePackage{etoolbox}

\newcommand{\ShorthandsOff}{\@ifundefined{shorthandoff}{}{\shorthandoff{;:!?}}}
\newcommand{\ShorthandsOn}{\@ifundefined{shorthandon}{}{\shorthandon{;:!?}}}
\ShorthandsOff
\AtBeginDocument{\ShorthandsOn}
\AtBeginEnvironment{tikzpicture}{\ShorthandsOff}

\RequirePackage{geometry}
\RequirePackage{fancyhdr}
\RequirePackage{lastpage}
\RequirePackage{multicol}
\RequirePackage{array}
\RequirePackage{enumitem}
\RequirePackage{xcolor}
\RequirePackage{needspace}

\RequirePackage{amsmath,amssymb}
\RequirePackage{mathpazo}
\DeclareMathSizes{10.95}{10.95}{8}{5.5}
\begingroup\catcode`\_=\active \catcode`\^=\active
\gdef_#1{\sb{\scriptscriptstyle #1}}
\gdef^#1{\sp{\scriptscriptstyle\lycee@moinsepais #1}}
\endgroup
\newcommand\lycee@moins{\mathbin{%
  \setbox\z@\hbox{$\m@th\scriptscriptstyle\lycee@moinsorig$}%
  \dimen@=.55\wd\z@ \advance\dimen@-.5pt
  \hbox to.75\wd\z@{\hss$\m@th\scriptscriptstyle\vcenter{\hbox{%
    \tikz\draw[line width=.5pt,line cap=round](0,0)--(\the\dimen@,0);}}$\hss}}}
\begingroup\lccode`\~=`\- \lowercase{\endgroup\let~}\lycee@moins
\newcommand\lycee@moinsepais{\mathcode`\-="8000 }
\AtBeginDocument{\mathchardef\lycee@moinsorig=\mathcode`\-\relax}
\AtBeginDocument{\catcode`\_=12 \mathcode`\_="8000
  \catcode`\^=12 \mathcode`\^="8000 }
\newcommand\lycee@exposants{%
  \fontdimen13\textfont2=.48\fontdimen6\textfont2
  \fontdimen14\textfont2=.48\fontdimen6\textfont2
  \fontdimen15\textfont2=.48\fontdimen6\textfont2
  \fontdimen13\scriptfont2=.48\fontdimen6\scriptfont2
  \fontdimen14\scriptfont2=.48\fontdimen6\scriptfont2
  \fontdimen15\scriptfont2=.48\fontdimen6\scriptfont2}
\every@math@size\expandafter{\the\every@math@size\lycee@exposants}
\newlength{\retraitcalculs}
\setlength{\retraitcalculs}{2em}

\RequirePackage{tikz}
\usetikzlibrary{arrows.meta,calc,babel,shapes.misc}
\RequirePackage[most]{tcolorbox}
\tcbuselibrary{skins,breakable}

\definecolor{coulDef}{HTML}{1F4E79}
\definecolor{coulProp}{HTML}{7030A0}
\definecolor{coulRem}{HTML}{C55A11}
\definecolor{coulEx}{HTML}{2E7D32}
\definecolor{coulExo}{HTML}{B03A2E}
\definecolor{coulPart}{HTML}{1F4E79}
\definecolor{coulMeth}{HTML}{1B6E4F}
\definecolor{coulCourbe}{HTML}{0B5ED7}
\definecolor{coulCourbeB}{HTML}{C00000}
\definecolor{coulCourbeC}{HTML}{2E7D32}
\definecolor{coulGrille}{HTML}{8C8C8C}
\definecolor{coulRep}{HTML}{1B6E4F}
\definecolor{coulBar}{HTML}{2E7D32}

\newcommand{\R}{\mathbb{R}}
\newcommand{\Rs}{\mathbb{R}^{*}}
\renewcommand{\le}{\leqslant}
\renewcommand{\ge}{\geqslant}
\renewcommand{\approx}{\simeq}
\newcommand{\vect}[1]{\overrightarrow{#1}}
\newcommand{\norme}[1]{\left\lVert #1 \right\rVert}
\newcommand{\intervalle}[2]{\left[#1\,;#2\right]}
\RequirePackage{cancel}
\renewcommand{\CancelColor}{\color{coulExo}}
\newcommand{\fois}[1]{\times{\color{coulExo}#1}}
\newcommand{\barre}[1]{\cancel{#1}}

\tikzset{grille/.style={coulGrille,line width=0.4pt}}
\newcommand{\quadrillage}[4]{\draw[grille,step=1] (#1,#2) grid (#3,#4);}
\tikzset{
  croix/.style={cross out,draw,line width=0.6pt,inner sep=0pt,minimum size=4pt},
  croixdroite/.style={croix,rotate=45,line width=0.8pt},
}
\newcommand{\pt}[3]{\path (#1) node[croix]{} node[#2,font=\small,inner sep=2.5pt]{$#3$};}

\newcommand{\lycee@repere}[4]{%
  \draw[grille,step=1] (#1,#2) grid (#3,#4);
  \draw[-{Stealth[length=2mm]},thick] (#1,0)--({#3+0.4},0) node[below left=-1pt]{$x$};
  \draw[-{Stealth[length=2mm]},thick] (0,#2)--(0,{#4+0.4}) node[below left=-1pt]{$y$};
  \node[below left,font=\scriptsize,inner sep=1.5pt] at (0,0) {$0$};}
\newcommand{\lycee@gradx}[1]{\foreach \k/\l in {#1}{%
  \draw (\k,0.07)--(\k,-0.07) node[below,font=\scriptsize,inner sep=1.5pt]{$\l$};}}
\newcommand{\lycee@grady}[1]{\foreach \k/\l in {#1}{%
  \draw (0.07,\k)--(-0.07,\k) node[left,font=\scriptsize,inner sep=1.5pt]{$\l$};}}
\newcommand{\lycee@intff}[2]{\left[#1\,;#2\right]}
\newcommand{\lycee@intfo}[2]{\left[#1\,;#2\right[}
\newcommand{\lycee@intof}[2]{\left]#1\,;#2\right]}
\newcommand{\lycee@intoo}[2]{\left]#1\,;#2\right[}
\newcommand{\lycee@ens}[1]{\left\{#1\right\}}
\AtBeginDocument{%
  \providecommand{\repere}{\lycee@repere}%
  \providecommand{\gradx}{\lycee@gradx}%
  \providecommand{\grady}{\lycee@grady}%
  \providecommand{\Cf}{\mathcal{C}\sb{\scriptscriptstyle f}}%
  \providecommand{\Cg}{\mathcal{C}\sb{\scriptscriptstyle g}}%
  \providecommand{\intff}{\lycee@intff}%
  \providecommand{\intfo}{\lycee@intfo}%
  \providecommand{\intof}{\lycee@intof}%
  \providecommand{\intoo}{\lycee@intoo}%
  \providecommand{\ens}{\lycee@ens}}

\newlist{questions}{enumerate}{3}
\setlist[questions,1]{label=\arabic*\textdegree),leftmargin=*,itemsep=2pt,topsep=3pt}
\setlist[questions,2]{label=\alph*),leftmargin=*,itemsep=1pt,topsep=2pt}
\setlist[questions,3]{label=\roman*),leftmargin=*,itemsep=1pt,topsep=1pt}
\setlist[itemize]{label=\textbullet,leftmargin=*,itemsep=2pt,topsep=3pt}

\newcommand{\besoinplace}[1]{\par\penalty\@M\begingroup
  \dimen@=#1\relax\dimen@ii=\pagegoal\advance\dimen@ii-\pagetotal
  \ifdim\dimen@>\dimen@ii\ifdim\dimen@ii>\z@\vfil\fi\break\fi\endgroup}

\newcommand{\lycee@pied}{\thepage/\pageref{LastPage}}
\newcommand{\entete}[2]{%
  \pagestyle{fancy}\fancyhf{}%
  \fancyhead[L]{\footnotesize\color{coulPart}\textbf{#1}}%
  \fancyhead[R]{\footnotesize\color{coulPart}\textbf{#2}}%
  \fancyfoot[C]{\footnotesize\lycee@pied}%
  \renewcommand{\headrulewidth}{0.4pt}%
  \renewcommand{\headrule}{\hbox to\headwidth{%
    \color{coulPart!40}\leaders\hrule height \headrulewidth\hfill}}}

\RequirePackage{ccicons}
\newcommand{\lycee@licence}{{\scriptsize\color{gray}%
  \copyright\ Luc Destombe --- \ccbyncsa\ CC BY-NC-SA 4.0}}
\newif\iflycee@licence \lycee@licencetrue
\newcommand{\sanslicence}{\lycee@licencefalse}
\AtBeginDocument{\iflycee@licence\AddToHookNext{shipout/foreground}{%
  \put(0,0){\rlap{\kern\dimexpr1in+\hoffset+\oddsidemargin+\textwidth\relax
    \llap{\raisebox{-\dimexpr1in+\voffset+\topmargin+\headheight+\headsep
      +\textheight+\footskip\relax}{\lycee@licence}}}}}\fi}

\newcommand{\formulesserrees}{%
  \setlength{\abovedisplayskip}{3pt}\setlength{\belowdisplayskip}{3pt}%
  \setlength{\abovedisplayshortskip}{2pt}\setlength{\belowdisplayshortskip}{3pt}}

\setlength{\parindent}{0pt}

\RequirePackage[implicit=false,hidelinks,pdfencoding=auto,psdextra,bookmarksopen,
  bookmarksopenlevel=1,pdfcreator={},pdfproducer={}]{hyperref}
\RequirePackage{bookmark}
\hypersetup{pdfauthor={Luc Destombe},pdfkeywords={CC BY-NC-SA 4.0}}
\newcounter{lycee@signet}
\newcommand{\signet}[3][]{\stepcounter{lycee@signet}%
  \edef\lycee@dest{\if\relax\detokenize{#1}\relax
    signet.\arabic{lycee@signet}\else#1\fi}%
  \hypertarget{\lycee@dest}{}\bookmark[level=#2,dest=\lycee@dest]{#3}}
\pdfstringdefDisableCommands{%
  \def\R{R}\def\vect#1{#1}\def\Cf{Cf}\def\Cg{Cg}\def\fois#1{×#1}%
  \def\barre#1{#1}\def\intervalle#1#2{[#1;#2]}\def\intff#1#2{[#1;#2]}%
  \def\textsuperscript#1{#1}\def\dfrac#1#2{#1/#2}\def\frac#1#2{#1/#2}%
  \def\vec#1{#1⃗}}%
\RequirePackage[official]{eurosym}

\tcbset{exercice corrige/.style={
  enhanced,breakable,before upper=\raggedright,
  colback=white,colframe=coulExo,
  boxrule=0.8pt,arc=2pt,
  left=6pt,right=6pt,top=8pt,bottom=5pt,
  before skip=9pt,after skip=4pt,
  coltitle=white,fonttitle=\bfseries\small,
  attach boxed title to top left={xshift=8pt,yshift=-\tcboxedtitleheight/2},
  boxed title style={colback=coulExo,arc=2pt,boxrule=0pt}}}
\newcounter{exo}
\newtcolorbox{exercice}[1][]{exercice corrige,
  phantom={\signet[exercice-\arabic{exo}]{2}{Exercice \arabic{exo}}},
  title={Exercice \theexo\ifx\relax#1\relax\else{} --- #1\fi}}
\newenvironment{exo}[1][\relax]{\refstepcounter{exo}\begin{exercice}[#1]}{\end{exercice}}

\geometry{top=1.8cm,bottom=1cm,left=1cm,right=1cm,headheight=14pt,footskip=0.6cm}
\renewcommand{\lycee@pied}{\thepage}

\renewcommand{\theexo}{\ifnum\value{exo}<10 0\fi\arabic{exo}}
\tcbset{exercice corrige/.style={
  enhanced,breakable,before upper=\raggedright,
  colback=white,colframe=coulExo,
  boxrule=0.9pt,arc=2pt,
  left=8pt,right=8pt,top=8pt,bottom=6pt,
  before skip=8pt,after skip=6pt,
  coltitle=white,fonttitle=\bfseries,
  attach boxed title to top left={xshift=8pt,yshift=-\tcboxedtitleheight/2},
  boxed title style={colback=coulExo,arc=2pt,boxrule=0pt}}}

\newcommand{\entetefiche}[2]{%
  \begin{center}
    \begin{tcolorbox}[enhanced,width=0.92\linewidth,colback=white,colframe=coulPart,
        boxrule=1.4pt,arc=3pt,halign=center,top=6pt,bottom=6pt]
      {\large\bfseries\color{coulPart} CORRIGÉ --- FICHE D'EXERCICES}\\[3pt]
      {\normalsize\bfseries #1}\\[2pt]
      {\normalsize #2}
    \end{tcolorbox}
  \end{center}}

\newcounter{partie}
\newcommand{\partie}[1]{%
  \stepcounter{partie}%
  \besoinplace{8\baselineskip}\vspace{4pt}%
  \begin{tcolorbox}[enhanced,colback=coulPart!8,colframe=coulPart,boxrule=0pt,
      leftrule=3pt,arc=0pt,left=6pt,right=4pt,top=3pt,bottom=3pt,
      before skip=6pt,after skip=2pt]
    \bfseries\color{coulPart}\Roman{partie} --- #1\signet{1}{\Roman{partie} --- #1}
  \end{tcolorbox}}
\newcommand{\partiesansnum}[1]{%
  \besoinplace{8\baselineskip}\vspace{4pt}%
  \begin{tcolorbox}[enhanced,colback=coulPart!8,colframe=coulPart,boxrule=0pt,
      leftrule=3pt,arc=0pt,left=6pt,right=4pt,top=3pt,bottom=3pt,
      before skip=6pt,after skip=2pt]
    \bfseries\color{coulPart}#1\signet{1}{#1}
  \end{tcolorbox}}

\setlist[questions,1]{label=\arabic*\textdegree),leftmargin=*,itemsep=3pt,topsep=3pt}
\setlist[questions,2]{label=\alph*),leftmargin=*,itemsep=2pt,topsep=2pt}
\setlist[questions,3]{label=\roman*),leftmargin=*,itemsep=1pt,topsep=1pt}
\newlist{expr}{enumerate}{1}
\setlist[expr]{label={\Alph*\ $=$},leftmargin=*,itemsep=2pt,topsep=3pt}

\newcommand{\res}[1]{{\color{coulRep}#1}}
\newcommand{\rem}[1]{\par\smallskip{\small\itshape\color{black!60}$\rightarrow$ #1}}
\newenvironment{calculs}{\par\smallskip\noindent$\displaystyle\begin{aligned}[t]}%
  {\end{aligned}$\par\smallskip}
\newcommand{\justif}[1]{\quad\text{\small\itshape\color{black!60}#1}}
\newcommand{\cas}[1]{\par\smallskip\noindent\textbf{\color{coulExo}#1}\nobreak}
\newcommand{\explique}[2]{\par\smallskip\noindent
  \begin{minipage}[t]{0.6\linewidth}\raggedright #1\end{minipage}\hfill
  \begin{minipage}[t]{0.37\linewidth}\raggedright\leavevmode
    {\small\itshape\color{black!60}#2}\end{minipage}\par}

\setlength{\parskip}{2pt}

\makeatother
\entete{Chapitre 1 --- Suites numériques --- corrigé}{1\textsuperscript{re} mathématiques spécifiques}
\setlength{\parindent}{0pt}
\setlength{\jot}{5pt}

\let\casligne\cas
\renewcommand{\cas}[1]{\casligne{#1}\par\nopagebreak}
\newcommand{\etiq}[1]{\textbf{\color{coulExo}#1}~}

\newcommand{\reperesuite}[5]{%
  \quadrillage{0}{#2}{#1}{#3}
  \draw[->,>=Stealth,black!70] (0,0) -- (#1+0.6,0) node[below,font=\footnotesize] {$n$};
  \draw[->,>=Stealth,black!70] (0,#2) -- (0,#3+0.6) node[left,font=\footnotesize] {$#5$};
  \foreach \i in {1,...,#1} {\node[below,font=\scriptsize,inner sep=1.5pt] at (\i,0) {$\i$};}
  \foreach \v in {#4} {\node[left,font=\scriptsize,inner sep=1.5pt] at (0,\v) {$\v$};}
  \node[below left,font=\scriptsize,inner sep=1.5pt] at (0,0) {$0$};}

\begin{document}

\entetefiche{Chapitre 1 --- Suites numériques}{1\textsuperscript{re} mathématiques spécifiques}

\partie{Notion de suite}

\begin{exo}[Calculer les premiers termes]
\explique{Les cinq premiers termes sont ceux de rang $0$ ; $1$ ; $2$ ; $3$ et $4$.}{le
premier terme est $u_{0}$, pas $u_{1}$}
\cas{$u_{n}=4n-5$}
\begin{calculs}
u_{0} &= 4\times 0-5 &&\justif{on remplace $n$ par $0$}\\
u_{0} &= \res{-5} &&\\
u_{1} &= 4\times 1-5 = \res{-1} &&\justif{on remplace $n$ par $1$}\\
u_{2} &= 4\times 2-5 = \res{3} &&\\
u_{3} &= 4\times 3-5 = \res{7} &&\\
u_{4} &= 4\times 4-5 = \res{11} &&\justif{d'un terme au suivant, on ajoute $4$}
\end{calculs}
\cas{$v_{n}=n^{2}+1$}
\begin{calculs}
v_{0} &= 0^{2}+1 = \res{1} &&\justif{$0^{2}=0\times 0=0$}\\
v_{1} &= 1^{2}+1 = \res{2} &&\justif{$1^{2}=1\times 1=1$}\\
v_{2} &= 2^{2}+1 = \res{5} &&\justif{le carré d'abord : $2^{2}=2\times 2=4$}\\
v_{3} &= 3^{2}+1 = \res{10} &&\justif{$3^{2}=9$}\\
v_{4} &= 4^{2}+1 = \res{17} &&\justif{$4^{2}=16$}
\end{calculs}
\cas{$w_{n}=3^{n}$}
\begin{calculs}
w_{0} &= 3^{0} = \res{1} &&\justif{un nombre à la puissance $0$ vaut $1$}\\
w_{1} &= 3^{1} = \res{3} &&\\
w_{2} &= 3^{2} = \res{9} &&\justif{$3\times 3$}\\
w_{3} &= 3^{3} = \res{27} &&\justif{$3\times 3\times 3$}\\
w_{4} &= 3^{4} = \res{81} &&\justif{$27\times 3$}
\end{calculs}
\rem{Erreurs fréquentes : $3^{0}=0$ (c'est $1$) ; $3^{2}=6$ (c'est $3\times 3=9$, pas
$3\times 2$).}
\end{exo}

\begin{exo}[Calculer un terme]
\cas{a)}
\begin{calculs}
u_{9} &= 6\times 9+5 &&\justif{on remplace $n$ par $9$}\\
u_{9} &= 54+5 &&\justif{la multiplication d'abord}\\
u_{9} &= \res{59} &&
\end{calculs}
\explique{On n'a pas besoin de $u_{0}$, $u_{1}$\ldots{} $u_{8}$.}{la formule donne
directement n'importe quel terme}
\cas{b)}
\begin{calculs}
v_{4} &= 3\times 4^{2}-2\times 4+1 &&\justif{on remplace chaque $n$ par $4$}\\
v_{4} &= 3\times 16-2\times 4+1 &&\justif{le carré d'abord : $4^{2}=16$}\\
v_{4} &= 48-8+1 &&\justif{puis les multiplications}\\
v_{4} &= \res{41} &&
\end{calculs}
\rem{Erreur fréquente : calculer $(3\times 4)^{2}=144$. Le carré ne porte que sur $n$ :
on calcule $4^{2}$, puis on multiplie par $3$.}
\cas{c)}
\begin{calculs}
w_{0} &= \dfrac{2\times 0}{0+1} = \dfrac{0}{1} &&\justif{on remplace $n$ par $0$, en
haut et en bas}\\
w_{0} &= \res{0} &&\\
w_{1} &= \dfrac{2\times 1}{1+1} = \dfrac{2}{2} &&\\
w_{1} &= \res{1} &&\\
w_{2} &= \dfrac{2\times 2}{2+1} = \res{\dfrac{4}{3}} &&\justif{la fraction ne se simplifie
pas}\\
w_{3} &= \dfrac{2\times 3}{3+1} = \dfrac{6}{4} &&\\
w_{3} &= \dfrac{3\times\barre{2}}{2\times\barre{2}} = \dfrac{3}{2} &&\justif{on simplifie
par $2$}\\
w_{3} &= \res{1{,}5} &&
\end{calculs}
\cas{d)}
\begin{calculs}
t_{3} &= 20-3^{2} &&\justif{on remplace $n$ par $3$}\\
t_{3} &= 20-9 &&\justif{$3^{2}=9$}\\
t_{3} &= \res{11} &&
\end{calculs}
\end{exo}

\begin{exo}[Des suites dans la vie courante]
\explique{Dans chaque situation, on cherche le montant de départ, puis ce qui s'ajoute
(ou se retire) à chaque entrée, minute ou mois : $u_{n}$ vaut le départ, plus (ou
moins) $n$ fois ce montant.}{on peut tester avec un petit nombre : $2$ entrées coûtent
$30+2\times 3{,}5$}
\cas{a)}
\begin{calculs}
u_{n} &= \res{30+3{,}5n} &&\justif{la carte, plus $n$ entrées à $3{,}50$~€}\\
u_{12} &= 30+3{,}5\times 12 &&\justif{on remplace $n$ par $12$}\\
u_{12} &= 30+42 &&\justif{$3{,}5\times 12=36+6$}\\
u_{12} &= \res{72} &&
\end{calculs}
\explique{Pour $12$ entrées dans l'année, on paie $72$~€, carte comprise.}{on répond
dans le contexte}
\cas{b)}
\begin{calculs}
u_{n} &= \res{1+0{,}25n} &&\justif{$1$~€ au départ, plus $0{,}25$~€ par minute}\\
u_{40} &= 1+0{,}25\times 40 &&\justif{on remplace $n$ par $40$}\\
u_{40} &= 1+10 &&\justif{$0{,}25\times 40$ : le quart de $40$}\\
u_{40} &= \res{11} &&
\end{calculs}
\explique{Un trajet de $40$ minutes coûte $11$~€.}{}
\cas{c)}
\begin{calculs}
u_{n} &= \res{9\,000-600n} &&\justif{la réserve baisse : on retire $600$ par mois}\\
u_{6} &= 9\,000-600\times 6 &&\justif{on remplace $n$ par $6$}\\
u_{6} &= 9\,000-3\,600 &&\\
u_{6} &= \res{5\,400} &&
\end{calculs}
\explique{Au bout de $6$ mois, la réserve est de $5\,400$~m$^{3}$.}{}
\rem{Une baisse se traduit par un signe $-$ devant le terme en $n$.}
\end{exo}

\begin{exo}[Ne pas confondre]
\cas{a)}
\begin{calculs}
u_{4} &= 2\times 4+5 = \res{13} &&\justif{on remplace $n$ par $4$}\\
u_{4}+1 &= 13+1 = \res{14} &&\justif{on ajoute $1$ au terme $u_{4}$}\\
u_{5} &= 2\times 5+5 = \res{15} &&\justif{le terme suivant : on remplace $n$ par $5$}
\end{calculs}
\explique{$u_{4}+1=14$, alors que $u_{5}=15$ : ce ne sont pas les mêmes nombres.}{$u_{n}+1$ :
le terme augmenté de $1$ ; $u_{n+1}$ : le terme suivant}
\cas{b)}
\begin{calculs}
u_{n+1} &= 2(n+1)+5 &&\justif{on remplace $n$ par $n+1$, entre parenthèses}\\
u_{n+1} &= 2n+2+5 &&\justif{on développe : $2\times n+2\times 1$}\\
u_{n+1} &= \res{2n+7} &&
\end{calculs}
\explique{Contrôle avec $n=4$ : $2\times 4+7=15$, c'est bien $u_{5}$.}{on retrouve la
question a)}
\rem{Erreur fréquente : écrire $u_{n+1}=2n+5+1$. C'est $u_{n}+1$, et non le terme
suivant.}
\end{exo}

\begin{exo}[Retrouver le rang]
\cas{a)}
\begin{calculs}
u_{0} &= 80-4\times 0 = \res{80} &&\justif{on remplace $n$ par $0$}\\
u_{1} &= 80-4\times 1 = \res{76} &&\\
u_{10} &= 80-4\times 10 = 80-40 &&\\
u_{10} &= \res{40} &&
\end{calculs}
\cas{b)}
\explique{On cherche $n$ tel que $u_{n}=48$ : on résout une équation dont l'inconnue est
le rang $n$.}{on connaît le terme, on cherche son rang}
\begin{calculs}
80-4n &= 48 &&\justif{on traduit $u_{n}=48$}\\
-4n &= -32 &&\justif{on retire $80$ des deux côtés}\\
n &= 8 &&\justif{on divise par $-4$ : $\dfrac{-32}{-4}=8$}
\end{calculs}
\explique{$u_{n}=48$ pour \res{$n=8$}.}{contrôle : $u_{8}=80-32=48$}
\begin{calculs}
80-4n &= 0 &&\justif{on traduit $u_{n}=0$}\\
-4n &= -80 &&\justif{on retire $80$}\\
n &= 20 &&\justif{on divise par $-4$}
\end{calculs}
\explique{$u_{n}=0$ pour \res{$n=20$}.}{contrôle : $u_{20}=80-80=0$}
\rem{Autre façon de voir : de $80$ à $48$, on a retiré $32$, c'est-à-dire $8$ fois $4$.}
\end{exo}

\partie{Suites définies par récurrence}

\begin{exo}[Calculer les termes pas à pas]
\explique{On part du premier terme et on applique la relation quatre fois : chaque terme
se calcule avec celui d'avant, dans l'ordre.}{suite définie par récurrence}
\cas{a)}
\begin{calculs}
u_{1} &= 3u_{0}-1 &&\justif{la relation avec $n=0$}\\
u_{1} &= 3\times 2-1 = \res{5} &&\justif{on remplace $u_{0}$ par $2$}\\
u_{2} &= 3\times 5-1 = \res{14} &&\justif{$u_{2}=3u_{1}-1$}\\
u_{3} &= 3\times 14-1 = \res{41} &&\justif{$u_{3}=3u_{2}-1$ ; $3\times 14=42$}\\
u_{4} &= 3\times 41-1 = \res{122} &&\justif{$3\times 41=123$}
\end{calculs}
\cas{b)}
\begin{calculs}
v_{1} &= 5+4 = \res{9} &&\justif{$v_{1}=v_{0}+4$ : on ajoute $4$}\\
v_{2} &= 9+4 = \res{13} &&\\
v_{3} &= 13+4 = \res{17} &&\\
v_{4} &= 17+4 = \res{21} &&
\end{calculs}
\cas{c)}
\begin{calculs}
w_{1} &= \dfrac{48}{2} = \res{24} &&\justif{$w_{1}=\dfrac{w_{0}}{2}$ : on prend la moitié}\\
w_{2} &= \dfrac{24}{2} = \res{12} &&\\
w_{3} &= \dfrac{12}{2} = \res{6} &&\\
w_{4} &= \dfrac{6}{2} = \res{3} &&
\end{calculs}
\cas{d)}
\begin{calculs}
t_{1} &= 20-6 = \res{14} &&\justif{$t_{1}=t_{0}-6$ : on retire $6$}\\
t_{2} &= 14-6 = \res{8} &&\\
t_{3} &= 8-6 = \res{2} &&\\
t_{4} &= 2-6 = \res{-4} &&\justif{on passe sous zéro}
\end{calculs}
\rem{On ne peut pas calculer $u_{4}$ directement : il faut $u_{3}$, donc $u_{2}$, donc
$u_{1}$.}
\end{exo}

\begin{exo}[Avec des nombres négatifs]
\cas{a)}
\begin{calculs}
u_{1} &= 3-u_{0}^{2} &&\justif{la relation avec $n=0$}\\
u_{1} &= 3-(-1)^{2} &&\justif{on remplace $u_{0}$ par $-1$, entre parenthèses}\\
u_{1} &= 3-1 &&\justif{$(-1)^{2}=(-1)\times(-1)=1$}\\
u_{1} &= \res{2} &&\\
u_{2} &= 3-2^{2} = 3-4 &&\justif{$u_{2}=3-u_{1}^{2}$}\\
u_{2} &= \res{-1} &&\\
u_{3} &= 3-(-1)^{2} = 3-1 &&\justif{même calcul que pour $u_{1}$}\\
u_{3} &= \res{2} &&
\end{calculs}
\explique{Les termes alternent : $2$ ; $-1$ ; $2$ ; $-1$\ldots}{dès qu'on retombe sur
$-1$, tout recommence}
\cas{b)}
\begin{calculs}
v_{1} &= -3v_{0}+2 &&\justif{la relation avec $n=0$}\\
v_{1} &= -3\times 1+2 = \res{-1} &&\justif{on remplace $v_{0}$ par $1$}\\
v_{2} &= -3\times(-1)+2 &&\justif{on remplace $v_{1}$ par $-1$, entre parenthèses}\\
v_{2} &= 3+2 = \res{5} &&\justif{moins par moins donne plus}\\
v_{3} &= -3\times 5+2 &&\\
v_{3} &= -15+2 = \res{-13} &&
\end{calculs}
\rem{Erreur fréquente : écrire $-1^{2}$ au lieu de $(-1)^{2}$. Sans parenthèses,
$-1^{2}$ se lit $-(1^{2})$, qui vaut $-1$. Un nombre négatif se remplace toujours entre
parenthèses.}
\end{exo}

\begin{exo}[Traduire un énoncé]
\explique{Le premier terme est la valeur de départ ($n=0$) ; la relation dit ce qu'on
fait pour passer d'une étape à la suivante.}{on cherche « ce qui se répète »}
\cas{a)}
\explique{$u_{0}=\res{150}$}{la somme au départ, après $0$ semaine}
\explique{$u_{n+1}=\res{u_{n}+20}$}{d'une semaine à la suivante, on ajoute $20$~€}
\begin{calculs}
u_{1} &= 150+20 = \res{170} &&\\
u_{2} &= 170+20 = \res{190} &&
\end{calculs}
\explique{Après $2$ semaines, la tirelire contient $190$~€.}{}
\cas{b)}
\explique{$a_{0}=\res{3}$}{la surface au départ}
\explique{$a_{n+1}=\res{3a_{n}}$}{« tripler », c'est multiplier par $3$}
\begin{calculs}
a_{1} &= 3\times 3 = \res{9} &&\\
a_{2} &= 3\times 9 = \res{27} &&
\end{calculs}
\explique{Après $2$ semaines, l'algue couvre $27$~m$^{2}$.}{}
\cas{c)}
\explique{$c_{0}=\res{400}$}{le nombre d'abonnés au départ}
\explique{Chaque année : $30$ départs et $50$ arrivées, soit $50-30=20$ abonnés de
plus.}{on fait le bilan de l'année}
\explique{$c_{n+1}=\res{c_{n}+20}$}{}
\begin{calculs}
c_{1} &= 400+20 = \res{420} &&\\
c_{2} &= 420+20 = \res{440} &&
\end{calculs}
\rem{Erreur fréquente : écrire $c_{n+1}=c_{n}-30$, en oubliant les arrivées. On regroupe
départs et arrivées en une seule variation.}
\end{exo}

\begin{exo}[Retrouver la relation]
\explique{On regarde comment passer d'un nombre au suivant : d'abord en ajoutant
(différences), puis en multipliant (quotients).}{on vérifie sur tous les nombres de la
liste}
\cas{a)}
\explique{$9-4=5$ ; \ $14-9=5$ ; \ $19-14=5$.}{on ajoute toujours $5$}
\explique{$u_{0}=\res{4}$ et $u_{n+1}=\res{u_{n}+5}$.}{}
\cas{b)}
\explique{$8-2=6$ mais $32-8=24$.}{ce n'est pas une addition}
\explique{$\dfrac{8}{2}=4$ ; \ $\dfrac{32}{8}=4$ ; \ $\dfrac{128}{32}=4$.}{on multiplie
toujours par $4$}
\explique{$u_{0}=\res{2}$ et $u_{n+1}=\res{4u_{n}}$.}{}
\cas{c)}
\explique{$4-1=3$ mais $13-4=9$ ; \ $\dfrac{4}{1}=4$ mais $4\times 4=16$, pas $13$.}{ni
une addition seule, ni une multiplication seule}
\explique{On essaie « multiplier par $3$, puis ajouter $1$ » : $3\times 1+1=4$ ; \
$3\times 4+1=13$ ; \ $3\times 13+1=40$.}{indice : les écarts $3$ ; $9$ ; $27$ sont
multipliés par $3$}
\explique{$u_{0}=\res{1}$ et $u_{n+1}=\res{3u_{n}+1}$.}{la relation marche pour toute la
liste}
\rem{Toujours vérifier la relation trouvée sur tous les termes de la liste, pas seulement
sur les deux premiers.}
\end{exo}

\begin{exo}[Revenir en arrière]
\explique{$u_{n+1}=u_{n}-6$ : pour avancer d'un rang, on retire $6$. Pour reculer d'un
rang, on fait l'inverse : on ajoute $6$.}{l'opération inverse de « $-6$ » est « $+6$ »}
\begin{calculs}
u_{3} &= u_{2}-6 &&\justif{la relation avec $n=2$}\\
11 &= u_{2}-6 &&\justif{on remplace $u_{3}$ par $11$}\\
u_{2} &= 11+6 = \res{17} &&\justif{on ajoute $6$ des deux côtés}\\
u_{1} &= 17+6 = \res{23} &&\justif{même raisonnement, un rang plus bas}\\
u_{0} &= 23+6 = \res{29} &&
\end{calculs}
\explique{Contrôle en avançant : $29$ ; $23$ ; $17$ ; $11$.}{on retire bien $6$ à chaque
fois}
\end{exo}

\begin{exo}[Au tableur]
\explique{Les rangs sont en colonne A ($0$ en A2, $1$ en A3\ldots), les termes en colonne
B (le premier terme en B2, le terme de rang $1$ en B3\ldots).}{la ligne 3 correspond au
rang $n=1$}
\cas{a)}
\explique{$u_{1}=u_{0}+15$ : le terme de la ligne 3 s'obtient à partir de celui de la
ligne 2, juste au-dessus.}{récurrence : on utilise la cellule du dessus}
\explique{En B3 : \res{\texttt{=B2+15}}}{recopiée vers le bas, elle devient \texttt{=B3+15}
en B4, et ainsi de suite}
\cas{b)}
\explique{$v_{1}=0{,}9\times v_{0}+30$, et $v_{0}$ est en B2.}{la multiplication s'écrit
\texttt{*}}
\explique{En B3 : \res{\texttt{=0,9*B2+30}}}{la virgule décimale reste une virgule}
\cas{c)}
\explique{$w_{n}=5n-1$ est une formule explicite : on a besoin du rang $n$, qui se trouve
en colonne A, sur la même ligne.}{en A3 : le rang $1$}
\explique{En B3 : \res{\texttt{=5*A3-1}}}{contrôle : $5\times 1-1=4$, c'est bien $w_{1}$}
\rem{Erreurs fréquentes : oublier le signe \texttt{=} au début ; écrire \texttt{=5n-1} (le
tableur ne connaît pas $n$, il connaît la cellule A3) ; écrire \texttt{5B2} au lieu de
\texttt{5*B2}.}
\end{exo}

\partie{Représentation graphique d'une suite}

\begin{exo}[Représenter une suite]
\cas{a)}
\begin{minipage}[c]{0.58\linewidth}\raggedright
\begin{calculs}
u_{0} &= 9-2\times 0 = 9 &&\justif{on remplace $n$ par $0$}\\
u_{1} &= 9-2\times 1 = 7 &&\justif{puis par $1$, et ainsi de suite}
\end{calculs}
\explique{Termes de rang $0$ à $6$ :}{on retire $2$ à chaque rang}
\smallskip
\renewcommand{\arraystretch}{1.2}
\begin{tabular}{|c|c|c|c|c|c|c|c|}
\hline
$n$ & $0$ & $1$ & $2$ & $3$ & $4$ & $5$ & $6$\\
\hline
$u_{n}$ & $9$ & $7$ & $5$ & $3$ & $1$ & $-1$ & $-3$\\
\hline
\end{tabular}

\smallskip
\explique{On place les points $(n\,;u_{n})$, sans les relier.}{$n$ en abscisse, $u_{n}$
en ordonnée}
\explique{Les points sont alignés.}{on descend de $2$ à chaque pas}
\end{minipage}\hfill
\begin{minipage}[c]{0.38\linewidth}\centering
\begin{tikzpicture}[x=0.4cm,y=0.4cm]
  \reperesuite{6}{-3}{9}{-2,2,4,6,8}{u_{n}}
  \foreach \n in {0,...,6} {\path[coulCourbe] (\n,9-2*\n) node[croix]{};}
\end{tikzpicture}
\end{minipage}
\cas{b)}
\begin{minipage}[c]{0.58\linewidth}\raggedright
\begin{calculs}
v_{1} &= 2-v_{0} = 2-4 &&\justif{la relation avec $n=0$}\\
v_{1} &= -2 &&\\
v_{2} &= 2-(-2) = 2+2 &&\justif{moins par moins donne plus}\\
v_{2} &= 4 &&\justif{on retrouve $v_{0}$ : tout recommence}
\end{calculs}
\renewcommand{\arraystretch}{1.2}
\begin{tabular}{|c|c|c|c|c|c|c|c|}
\hline
$n$ & $0$ & $1$ & $2$ & $3$ & $4$ & $5$ & $6$\\
\hline
$v_{n}$ & $4$ & $-2$ & $4$ & $-2$ & $4$ & $-2$ & $4$\\
\hline
\end{tabular}

\smallskip
\explique{On place les points $(n\,;v_{n})$, sans les relier.}{les termes alternent entre
$4$ et $-2$}
\end{minipage}\hfill
\begin{minipage}[c]{0.38\linewidth}\centering
\begin{tikzpicture}[x=0.45cm,y=0.45cm]
  \reperesuite{6}{-2}{4}{-2,-1,1,2,3,4}{v_{n}}
  \foreach \n/\v in {0/4,1/-2,2/4,3/-2,4/4,5/-2,6/4}
    {\path[coulCourbe] (\n,\v) node[croix]{};}
\end{tikzpicture}
\end{minipage}
\rem{On ne relie jamais les points : $n$ ne prend que des valeurs entières.}
\end{exo}

\begin{exo}[Lire un graphique]
\cas{a)}
\explique{$u_{0}=\res{1}$}{le point d'abscisse $0$ a pour ordonnée $1$}
\explique{$u_{4}=\res{5{,}5}$}{le point d'abscisse $4$ est au milieu, entre les lignes $5$
et $6$}
\cas{b)}
\explique{Le point d'ordonnée $4$ a pour abscisse $3$ : $u_{n}=4$ pour \res{$n=3$}.}{on
lit dans l'autre sens : de l'ordonnée vers l'abscisse}
\cas{c)}
\explique{Chaque point est plus haut que le précédent : les termes \res{augmentent}.}{on
lit de gauche à droite}
\rem{Ne pas confondre le rang $n$ (abscisse, en bas) et le terme $u_{n}$ (ordonnée, à
gauche).}
\end{exo}

\begin{exo}[Calculer puis représenter]
\cas{a)}
\begin{calculs}
u_{1} &= \dfrac{u_{0}+10}{2} &&\justif{la relation avec $n=0$}\\
u_{1} &= \dfrac{-6+10}{2} = \dfrac{4}{2} &&\justif{on remplace $u_{0}$ par $-6$}\\
u_{1} &= \res{2} &&\\
u_{2} &= \dfrac{2+10}{2} = \dfrac{12}{2} &&\justif{on ajoute $10$, puis on prend la
moitié}\\
u_{2} &= \res{6} &&\\
u_{3} &= \dfrac{6+10}{2} = \dfrac{16}{2} &&\\
u_{3} &= \res{8} &&\\
u_{4} &= \dfrac{8+10}{2} = \dfrac{18}{2} &&\\
u_{4} &= \res{9} &&
\end{calculs}
\cas{b)}
\begin{minipage}[c]{0.58\linewidth}\raggedright
\renewcommand{\arraystretch}{1.2}
\begin{tabular}{|c|c|c|c|c|c|}
\hline
$n$ & $0$ & $1$ & $2$ & $3$ & $4$\\
\hline
$u_{n}$ & $-6$ & $2$ & $6$ & $8$ & $9$\\
\hline
\end{tabular}

\smallskip
\explique{On place les points $(n\,;u_{n})$, sans les relier.}{le premier point,
$(0\,;-6)$, est sous l'axe des abscisses}
\explique{Les termes augmentent, de moins en moins vite.}{écarts : $8$ ; $4$ ; $2$ ; $1$}
\end{minipage}\hfill
\begin{minipage}[c]{0.38\linewidth}\centering
\begin{tikzpicture}[x=0.32cm,y=0.32cm]
  \reperesuite{4}{-6}{9}{-6,-4,-2,2,4,6,8}{u_{n}}
  \foreach \n/\v in {0/-6,1/2,2/6,3/8,4/9} {\path[coulCourbe] (\n,\v) node[croix]{};}
\end{tikzpicture}
\end{minipage}
\end{exo}

\begin{exo}[Une suite qui descend puis remonte]
\cas{a)}
\begin{calculs}
w_{0} &= 0^{2}-6\times 0+8 = 8 &&\justif{on remplace chaque $n$ par $0$}\\
w_{1} &= 1^{2}-6\times 1+8 = 1-6+8 &&\\
w_{1} &= 3 &&\\
w_{2} &= 2^{2}-6\times 2+8 = 4-12+8 &&\justif{le carré, puis la multiplication}\\
w_{2} &= 0 &&\\
w_{3} &= 3^{2}-6\times 3+8 = 9-18+8 &&\\
w_{3} &= -1 &&\\
w_{4} &= 4^{2}-6\times 4+8 = 16-24+8 &&\\
w_{4} &= 0 &&\\
w_{5} &= 5^{2}-6\times 5+8 = 25-30+8 &&\\
w_{5} &= 3 &&\\
w_{6} &= 6^{2}-6\times 6+8 = 36-36+8 &&\\
w_{6} &= 8 &&
\end{calculs}
\renewcommand{\arraystretch}{1.2}
\begin{tabular}{|c|c|c|c|c|c|c|c|}
\hline
$n$ & $0$ & $1$ & $2$ & $3$ & $4$ & $5$ & $6$\\
\hline
$w_{n}$ & \res{$8$} & \res{$3$} & \res{$0$} & \res{$-1$} & \res{$0$} & \res{$3$} &
\res{$8$}\\
\hline
\end{tabular}
\cas{b)}
\begin{minipage}[c]{0.58\linewidth}\raggedright
\explique{On place les points $(n\,;w_{n})$ du tableau, sans les relier.}{$n$ en
abscisse, $w_{n}$ en ordonnée}
\end{minipage}\hfill
\begin{minipage}[c]{0.38\linewidth}\centering
\begin{tikzpicture}[x=0.4cm,y=0.4cm]
  \reperesuite{6}{-1}{8}{2,4,6,8}{w_{n}}
  \foreach \n/\v in {0/8,1/3,2/0,3/-1,4/0,5/3,6/8} {\path[coulCourbe] (\n,\v) node[croix]{};}
\end{tikzpicture}
\end{minipage}
\cas{c)}
\explique{Les termes \res{diminuent} de $w_{0}=8$ jusqu'à $w_{3}=-1$, puis
\res{augmentent} jusqu'à $w_{6}=8$.}{le point le plus bas est $(3\,;-1)$}
\end{exo}

\partie{Suites arithmétiques : définition}

\begin{exo}[Reconnaître une suite arithmétique]
\explique{Pour chaque liste, on calcule les différences « terme suivant moins terme
précédent ». Si elles sont toutes égales, la suite est arithmétique, et la raison est
cette différence.}{propriété 1}
\cas{a)}
\explique{$10-6=4$ ; \ $14-10=4$ ; \ $18-14=4$.}{différences successives}
\explique{Arithmétique, de raison \res{$r=4$}.}{les différences sont égales}
\cas{b)}
\explique{$9-12=-3$ ; \ $6-9=-3$ ; \ $3-6=-3$.}{}
\explique{Arithmétique, de raison \res{$r=-3$}.}{les termes diminuent : la raison est
négative}
\cas{c)}
\explique{$6-2=4$ mais $18-6=12$.}{}
\explique{\res{Pas arithmétique}.}{les différences changent ; ici, on multiplie par $3$}
\cas{d)}
\explique{$2-1{,}5=0{,}5$ ; \ $2{,}5-2=0{,}5$ ; \ $3-2{,}5=0{,}5$.}{}
\explique{Arithmétique, de raison \res{$r=0{,}5$}.}{une raison peut être décimale}
\cas{e)}
\explique{$1-0=1$ mais $4-1=3$.}{}
\explique{\res{Pas arithmétique}.}{ce sont les carrés $0^{2}$ ; $1^{2}$ ; $2^{2}$ ;
$3^{2}$}
\rem{Erreur fréquente : calculer « précédent moins suivant » ($12-9=3$) et trouver une
raison positive pour une suite qui diminue.}
\end{exo}

\begin{exo}[Modéliser]
\explique{Dans chaque situation, on ajoute ou on retire la même quantité à chaque étape :
la raison est cette quantité, positive pour une hausse, négative pour une baisse.}{même
quantité à chaque étape : suite arithmétique}
\cas{a)}
\explique{$u_{0}=\res{50\,000}$}{les places vendues la première année}
\explique{$u_{n+1}=\res{u_{n}+4\,000}$}{chaque année, $4\,000$ places de plus}
\explique{Raison : \res{$r=4\,000$}.}{}
\cas{b)}
\explique{On note $u_{n}$ le stock, en cartons, après $n$ semaines.}{on nomme la suite}
\explique{$u_{0}=\res{6\,000}$}{le stock au départ}
\explique{$u_{n+1}=\res{u_{n}-450}$}{chaque semaine, $450$ cartons de moins}
\explique{Raison : \res{$r=-450$}.}{retirer $450$, c'est ajouter $-450$}
\cas{c)}
\explique{On note $u_{n}$ la somme, en euros, après $n$ mois.}{}
\explique{$u_{0}=\res{200}$}{la somme au départ}
\explique{$u_{n+1}=\res{u_{n}+15}$}{chaque mois, $15$~€ de plus}
\explique{Raison : \res{$r=15$}.}{}
\end{exo}

\begin{exo}[Arithmétique ou non ?]
\cas{a)}
\explique{On calcule d'abord $u_{n+1}$ : on remplace $n$ par $n+1$.}{comme à
l'exercice 4}
\begin{calculs}
u_{n+1} &= 7-3(n+1) &&\justif{$n+1$ entre parenthèses}\\
u_{n+1} &= 7-3n-3 &&\justif{on développe : $-3\times n$ et $-3\times 1$}\\
u_{n+1} &= 4-3n &&\justif{$7-3=4$}\\
u_{n+1}-u_{n} &= (4-3n)-(7-3n) &&\justif{on retire tout $u_{n}$ : parenthèses}\\
u_{n+1}-u_{n} &= 4-3n-7+3n &&\justif{le $-$ devant la parenthèse change les signes}\\
u_{n+1}-u_{n} &= -3 &&\justif{$-3n+3n=0$ et $4-7=-3$}
\end{calculs}
\explique{La différence vaut toujours $-3$ : la suite est \res{arithmétique, de raison
$-3$}.}{la différence ne dépend pas de $n$}
\cas{b)}
\begin{calculs}
v_{0} &= 0^{2}+3 = 3 &&\\
v_{1} &= 1^{2}+3 = 4 &&\\
v_{2} &= 2^{2}+3 = 7 &&
\end{calculs}
\explique{$v_{1}-v_{0}=1$ mais $v_{2}-v_{1}=3$ : la suite \res{n'est pas
arithmétique}.}{un seul contre-exemple suffit}
\cas{c)}
\begin{calculs}
w_{n} &= 3(n+2) &&\justif{on recopie la formule}\\
w_{n} &= 3n+6 &&\justif{on développe}\\
w_{n+1} &= 3(n+1)+6 &&\justif{on remplace $n$ par $n+1$}\\
w_{n+1} &= 3n+3+6 &&\\
w_{n+1} &= 3n+9 &&\\
w_{n+1}-w_{n} &= (3n+9)-(3n+6) &&\\
w_{n+1}-w_{n} &= 3n+9-3n-6 &&\justif{le $-$ change les signes}\\
w_{n+1}-w_{n} &= 3 &&
\end{calculs}
\explique{La différence vaut toujours $3$ : la suite est \res{arithmétique, de raison
$3$}.}{premier terme : $w_{0}=3\times 2=6$}
\rem{Pour montrer qu'une suite est arithmétique, quelques termes ne suffisent pas : on
calcule $u_{n+1}-u_{n}$ avec la lettre $n$. Pour montrer qu'elle ne l'est pas, un seul
contre-exemple suffit.}
\end{exo}

\begin{exo}[Compléter]
\cas{a)}
\begin{calculs}
u_{1} &= u_{0}+r &&\justif{arithmétique : on ajoute la raison}\\
u_{1} &= 9+(-4) = \res{5} &&\justif{ajouter $-4$, c'est retirer $4$}\\
u_{2} &= 5-4 = \res{1} &&\\
u_{3} &= 1-4 = \res{-3} &&\justif{on passe sous zéro}\\
u_{4} &= -3-4 = \res{-7} &&
\end{calculs}
\cas{b)}
\begin{calculs}
r &= v_{3}-v_{2} &&\justif{différence de deux termes consécutifs}\\
r &= 18-13 = \res{5} &&\\
v_{4} &= v_{3}+5 = 18+5 &&\justif{on avance d'un rang : on ajoute $5$}\\
v_{4} &= \res{23} &&\\
v_{1} &= v_{2}-5 = 13-5 &&\justif{on recule d'un rang : on retire $5$}\\
v_{1} &= \res{8} &&\\
v_{0} &= v_{1}-5 = 8-5 &&\\
v_{0} &= \res{3} &&
\end{calculs}
\explique{Contrôle : $3$ ; $8$ ; $13$ ; $18$ ; $23$.}{on ajoute bien $5$ à chaque fois}
\end{exo}

\partie{Suites arithmétiques : expression en fonction de $n$}

\begin{exo}[Appliquer la formule]
\explique{On utilise $u_{n}=u_{0}+n\times r$ : il faut le premier terme $u_{0}$ et la
raison $r$.}{propriété 2}
\cas{a)}
\begin{calculs}
u_{n} &= u_{0}+n\times r &&\justif{premier terme $10$, raison $-3$}\\
u_{n} &= 10+n\times(-3) &&\\
u_{n} &= \res{10-3n} &&\justif{ajouter $-3n$, c'est retirer $3n$}\\
u_{8} &= 10-3\times 8 &&\justif{on remplace $n$ par $8$}\\
u_{8} &= 10-24 &&\\
u_{8} &= \res{-14} &&
\end{calculs}
\cas{b)}
\begin{calculs}
u_{n} &= u_{0}+n\times r &&\justif{premier terme $-4$, raison $5$}\\
u_{n} &= \res{-4+5n} &&\\
u_{12} &= -4+5\times 12 &&\justif{on remplace $n$ par $12$}\\
u_{12} &= -4+60 &&\\
u_{12} &= \res{56} &&
\end{calculs}
\cas{c)}
\explique{Le premier terme donné est $u_{1}$, pas $u_{0}$ : on recule d'un rang.}{la
formule part de $u_{0}$}
\begin{calculs}
u_{0} &= u_{1}-r &&\justif{on recule d'un rang : on retire la raison}\\
u_{0} &= 600-40 = 560 &&\\
u_{n} &= \res{560+40n} &&\justif{$u_{0}+n\times r$}\\
u_{8} &= 560+40\times 8 &&\justif{on remplace $n$ par $8$}\\
u_{8} &= 560+320 &&\\
u_{8} &= \res{880} &&
\end{calculs}
\explique{Contrôle : $u_{1}=560+40=600$.}{on retrouve la donnée}
\cas{d)}
\begin{calculs}
u_{0} &= u_{1}-r &&\justif{on recule d'un rang}\\
u_{0} &= -7-4 = -11 &&\\
u_{n} &= \res{-11+4n} &&\justif{$u_{0}+n\times r$}\\
u_{5} &= -11+4\times 5 &&\\
u_{5} &= -11+20 &&\\
u_{5} &= \res{9} &&
\end{calculs}
\explique{Contrôle : $u_{1}=-11+4=-7$.}{on retrouve la donnée}
\rem{Erreur fréquente : mettre $u_{1}$ à la place de $u_{0}$ dans la formule. En c), on
trouverait $600+40\times 8=920$, ce qui est faux.}
\end{exo}

\begin{exo}[Une suite qui diminue]
\cas{a)}
\explique{$172-185=-13$ ; \ $159-172=-13$ ; \ $146-159=-13$.}{différences successives}
\explique{Les différences sont égales : ce sont les termes d'une suite arithmétique, de
raison \res{$r=-13$}.}{propriété 1}
\cas{b)}
\begin{calculs}
u_{5} &= u_{4}+r &&\justif{on avance d'un rang}\\
u_{5} &= 146-13 = \res{133} &&
\end{calculs}
\explique{Pour $u_{14}$, on ne calcule pas tous les termes : on exprime $u_{n}$ en
fonction de $n$.}{il faudrait neuf calculs de plus}
\begin{calculs}
u_{0} &= u_{1}-r &&\justif{on recule d'un rang}\\
u_{0} &= 185-(-13) = 185+13 &&\justif{retirer $-13$, c'est ajouter $13$}\\
u_{0} &= 198 &&\\
u_{n} &= 198-13n &&\justif{$u_{0}+n\times r$}\\
u_{14} &= 198-13\times 14 &&\justif{on remplace $n$ par $14$}\\
u_{14} &= 198-182 &&\justif{$13\times 14=130+52$}\\
u_{14} &= \res{16} &&
\end{calculs}
\explique{Contrôle : $u_{1}=198-13=185$.}{on retrouve la donnée}
\end{exo}

\begin{exo}[Production d'un atelier]
\cas{a)}
\explique{Chaque année, la production augmente de la même quantité, $250$ vélos : la suite
$(u_{n})$ est \res{arithmétique}, de premier terme $u_{0}=4\,000$ et de raison
$r=250$.}{$u_{0}$ : la production de 2010, car $2010+0=2010$}
\begin{calculs}
u_{n} &= u_{0}+n\times r &&\\
u_{n} &= \res{4\,000+250n} &&
\end{calculs}
\cas{b)}
\begin{calculs}
u_{3} &= 4\,000+250\times 3 &&\justif{2013 correspond à $n=3$, car $2010+3=2013$}\\
u_{3} &= 4\,000+750 &&\\
u_{3} &= \res{4\,750} &&
\end{calculs}
\explique{En 2013, l'atelier fabrique $4\,750$ vélos.}{}
\begin{calculs}
u_{20} &= 4\,000+250\times 20 &&\justif{2030 correspond à $n=20$}\\
u_{20} &= 4\,000+5\,000 &&\justif{$250\times 20=250\times 2\times 10$}\\
u_{20} &= \res{9\,000} &&
\end{calculs}
\explique{En 2030, l'atelier fabriquera $9\,000$ vélos.}{}
\rem{Erreur fréquente : remplacer $n$ par l'année ($2013$) au lieu du nombre d'années
écoulées depuis 2010 ($3$).}
\end{exo}

\begin{exo}[Lire la formule]
\explique{Une suite arithmétique s'écrit $u_{n}=u_{0}+r\times n$ : le nombre seul est le
premier terme, le nombre qui multiplie $n$ est la raison.}{on compare avec la formule}
\cas{a)}
\explique{$u_{0}=\res{5}$ et $r=\res{4}$.}{contrôle : $u_{0}=5+4\times 0=5$}
\cas{b)}
\explique{$v_{0}=\res{18}$ et $r=\res{-1{,}5}$.}{le signe $-$ fait partie de la raison}
\cas{c)}
\explique{$w_{n}=6n$ s'écrit $w_{n}=0+6n$ : $w_{0}=\res{0}$ et $r=\res{6}$.}{il n'y a pas
de nombre seul : le premier terme vaut $0$}
\end{exo}

\begin{exo}[Un terme lointain]
\cas{a)}
\begin{calculs}
u_{n} &= 15+4n &&\justif{$u_{0}+n\times r$ : premier terme $15$, raison $4$}\\
u_{100} &= 15+4\times 100 &&\justif{on remplace $n$ par $100$}\\
u_{100} &= 15+400 &&\\
u_{100} &= \res{415} &&
\end{calculs}
\cas{b)}
\begin{calculs}
u_{0} &= u_{1}-r &&\justif{le premier terme donné est $u_{1}$ : on recule d'un rang}\\
u_{0} &= 7-3 = 4 &&\\
u_{n} &= 4+3n &&\justif{$u_{0}+n\times r$}\\
u_{50} &= 4+3\times 50 &&\\
u_{50} &= 4+150 &&\\
u_{50} &= \res{154} &&
\end{calculs}
\rem{Écrire $u_{50}=7+50\times 3=157$ est faux : de $u_{1}$ à $u_{50}$, on n'ajoute la
raison que $49$ fois. Passer par $u_{0}$ évite le piège.}
\end{exo}

\partie{Suites arithmétiques : problème de seuil}

\begin{exo}[Acheter un vélo électrique]
\cas{a)}
\begin{calculs}
u_{2} &= 900+60 = \res{960} &&\justif{le 1\textsuperscript{er} février : on ajoute $60$}\\
u_{3} &= 960+60 = \res{1\,020} &&\justif{le 1\textsuperscript{er} mars}\\
u_{4} &= 1\,020+60 = \res{1\,080} &&\justif{le 1\textsuperscript{er} avril}
\end{calculs}
\cas{b)}
\explique{Elle ajoute $60$~€ chaque mois : $(u_{n})$ est arithmétique, de raison
$r=60$.}{même somme ajoutée chaque mois}
\explique{Le premier terme donné est $u_{1}$ : on recule d'un rang pour trouver
$u_{0}$.}{la formule part de $u_{0}$}
\begin{calculs}
u_{0} &= u_{1}-r &&\justif{on retire la raison}\\
u_{0} &= 900-60 = 840 &&\\
u_{n} &= u_{0}+n\times r &&\\
u_{n} &= \res{840+60n} &&
\end{calculs}
\explique{Contrôle : $u_{1}=840+60=900$.}{on retrouve la donnée}
\cas{c)}
\begin{calculs}
u_{7} &= 840+60\times 7 &&\justif{on remplace $n$ par $7$}\\
u_{7} &= 840+420 &&\\
u_{7} &= \res{1\,260} &&
\end{calculs}
\explique{Le 1\textsuperscript{er} juillet, elle aura $1\,260$~€.}{$u_{7}$ : le
septième mois, juillet}
\cas{d)}
\begin{calculs}
840+60n &\ge 1\,500 &&\justif{on traduit : au moins $1\,500$~€}\\
60n &\ge 660 &&\justif{on retire $840$}\\
n &\ge 11 &&\justif{on divise par $60$, positif : $\dfrac{660}{60}=11$}
\end{calculs}
\explique{Elle aura $1\,500$~€ pour \res{$n=11$}.}{contrôle : $u_{11}=840+660=1\,500$}
\cas{e)}
\explique{$u_{1}$ est la somme au 1\textsuperscript{er} janvier, $u_{2}$ au
1\textsuperscript{er} février\ldots{} $u_{11}$ au 1\textsuperscript{er} du onzième mois,
novembre.}{le rang $n$ est le numéro du mois}
\explique{Elle pourra acheter son vélo \res{le 1\textsuperscript{er} novembre}.}{on
répond par une date, comme demandé}
\end{exo}

\begin{exo}[Doubler la production]
\explique{En 2010, la production est $u_{0}=4\,000$ ; « doubler » veut dire atteindre
$2\times 4\,000=8\,000$ vélos.}{on traduit la question}
\begin{calculs}
4\,000+250n &\ge 8\,000 &&\justif{au moins le double}\\
250n &\ge 4\,000 &&\justif{on retire $4\,000$}\\
n &\ge 16 &&\justif{on divise par $250$, positif : $250\times 16=4\,000$}
\end{calculs}
\explique{Le premier entier qui convient est $n=16$ : la production aura doublé \res{en
2026}.}{$2010+16=2026$ ; contrôle : $u_{16}=4\,000+4\,000=8\,000$}
\end{exo}

\begin{exo}[Le recul d'une falaise]
\cas{a)}
\explique{La distance baisse de $2$~m par an : $(d_{n})$ est arithmétique, de premier
terme $d_{0}=200$ et de raison $r=-2$.}{$d_{0}$ : la distance en 2020}
\begin{calculs}
d_{n} &= d_{0}+n\times r &&\\
d_{n} &= 200+n\times(-2) &&\\
d_{n} &= \res{200-2n} &&
\end{calculs}
\cas{b)}
\begin{calculs}
200-2n &< 175 &&\justif{on traduit : « à moins de $175$~m »}\\
-2n &< -25 &&\justif{on retire $200$}\\
n &> 12{,}5 &&\justif{on divise par $-2$, négatif : le sens change}
\end{calculs}
\explique{Le premier entier qui convient est $n=13$ : la maison sera à moins de $175$~m
du bord à partir de \res{2033}.}{$2020+13=2033$ ; contrôle : $d_{12}=176$ et
$d_{13}=174$}
\rem{Erreur fréquente : oublier de changer le sens en divisant par $-2$. On trouverait
$n<12{,}5$, absurde ici : la distance diminue, elle finit par passer sous $175$~m.}
\end{exo}

\begin{exo}[Résoudre des inéquations]
\explique{On résout comme une équation, puis on cherche le plus petit entier qui
convient.}{diviser par un nombre négatif change le sens}
\cas{a)}
\begin{calculs}
15+5n &> 100 &&\justif{on recopie l'inéquation}\\
5n &> 85 &&\justif{on retire $15$}\\
n &> 17 &&\justif{on divise par $5$, positif}
\end{calculs}
\explique{$17$ ne convient pas : le plus petit entier est \res{$n=18$}.}{« $>$ » :
$15+5\times 17=100$, qui n'est pas plus grand que $100$}
\cas{b)}
\begin{calculs}
240-30n &< 0 &&\\
-30n &< -240 &&\justif{on retire $240$}\\
n &> 8 &&\justif{on divise par $-30$, négatif : le sens change}
\end{calculs}
\explique{Le plus petit entier est \res{$n=9$}.}{$n=8$ donne exactement $0$, qui n'est pas
négatif}
\cas{c)}
\begin{calculs}
40+6n &\ge 100 &&\\
6n &\ge 60 &&\justif{on retire $40$}\\
n &\ge 10 &&\justif{on divise par $6$, positif}
\end{calculs}
\explique{Le plus petit entier est \res{$n=10$}.}{« $\ge$ » : $10$ convient, car
$40+60=100$}
\rem{Bien regarder le symbole : avec $>$ ou $<$, la valeur limite est exclue ; avec $\ge$
ou $\le$, elle est comprise.}
\end{exo}

\partie{Suites arithmétiques : variations et représentation graphique}

\begin{exo}[Sens de variation]
\explique{Il suffit de la raison : si $r>0$, la suite est croissante ; si $r<0$, elle est
décroissante.}{propriété 3}
\cas{a)}
\explique{$u_{n}=3+8n$ : la raison est le nombre qui multiplie $n$, $r=\res{8}$.}{forme
$u_{0}+r\times n$}
\explique{$8>0$ : $(u_{n})$ est \res{croissante}.}{}
\cas{b)}
\explique{$v_{n}=20-0{,}4n$ : $r=\res{-0{,}4}$.}{le signe $-$ fait partie de la raison}
\explique{$-0{,}4<0$ : $(v_{n})$ est \res{décroissante}.}{}
\cas{c)}
\explique{$t_{n}=-5n$ s'écrit $0-5n$ : $r=\res{-5}$.}{premier terme $0$}
\explique{$-5<0$ : $(t_{n})$ est \res{décroissante}.}{}
\cas{d)}
\explique{$w_{n+1}=w_{n}+0{,}5$ : on ajoute $0{,}5$ à chaque étape, $r=\res{0{,}5}$.}{}
\explique{$0{,}5>0$ : $(w_{n})$ est \res{croissante}.}{}
\rem{Le premier terme, même négatif ($w_{0}=-6$), ne compte pas : seul le signe de la
raison décide.}
\end{exo}

\begin{exo}[Lire un graphique]
\cas{a)}
\explique{$u_{0}=\res{8}$}{le premier point, d'abscisse $0$, a pour ordonnée $8$}
\explique{D'un point au suivant, on descend de $1{,}5$ : de $8$ à $6{,}5$, puis à $5$\ldots{}
Donc \res{$r=-1{,}5$}.}{on descend : la raison est négative}
\cas{b)}
\begin{calculs}
u_{n} &= u_{0}+n\times r &&\justif{premier terme $8$, raison $-1{,}5$}\\
u_{n} &= 8+n\times(-1{,}5) &&\\
u_{n} &= \res{8-1{,}5n} &&
\end{calculs}
\cas{c)}
\begin{calculs}
u_{10} &= 8-1{,}5\times 10 &&\justif{on remplace $n$ par $10$}\\
u_{10} &= 8-15 &&\\
u_{10} &= \res{-7} &&
\end{calculs}
\cas{d)}
\explique{$r=-1{,}5<0$ : la suite est \res{décroissante}.}{les points descendent}
\end{exo}

\begin{exo}[Représenter]
\begin{minipage}[c]{0.58\linewidth}\raggedright
\cas{a)}
\explique{$u_{0}=1$ ; puis on ajoute $1{,}5$ à chaque rang.}{$u_{n}=1+1{,}5n$ : raison
$1{,}5$}
\renewcommand{\arraystretch}{1.2}
\begin{tabular}{|c|c|c|c|c|c|}
\hline
$n$ & $0$ & $1$ & $2$ & $3$ & $4$\\
\hline
$u_{n}$ & $1$ & $2{,}5$ & $4$ & $5{,}5$ & $7$\\
\hline
\end{tabular}

\explique{$r=1{,}5>0$ : $(u_{n})$ est \res{croissante}.}{croix $\times$ : les points
montent}
\cas{b)}
\explique{$v_{0}=8$ ; puis on retire $2$ à chaque rang.}{raison $-2$}
\renewcommand{\arraystretch}{1.2}
\begin{tabular}{|c|c|c|c|c|c|}
\hline
$n$ & $0$ & $1$ & $2$ & $3$ & $4$\\
\hline
$v_{n}$ & $8$ & $6$ & $4$ & $2$ & $0$\\
\hline
\end{tabular}

\explique{$r=-2<0$ : $(v_{n})$ est \res{décroissante}.}{croix $+$ : les points
descendent}
\end{minipage}\hfill
\begin{minipage}[c]{0.38\linewidth}\centering
\begin{tikzpicture}[x=0.5cm,y=0.5cm]
  \reperesuite{4}{0}{8}{1,2,3,4,5,6,7,8}{}
  \foreach \n/\v in {0/1,1/2.5,2/4,3/5.5,4/7} {\path[coulCourbe] (\n,\v) node[croix]{};}
  \foreach \n/\v in {0/8,1/6,2/4,3/2,4/0} {\path[coulCourbeB] (\n,\v) node[croixdroite]{};}
\end{tikzpicture}

{\small $\times$ : $(u_{n})$ \quad $+$ : $(v_{n})$}
\end{minipage}
\rem{Les points de chaque suite sont alignés : ce sont des suites arithmétiques. Au rang
$2$, les deux suites valent $4$ : les points sont confondus.}
\end{exo}

\partie{Suites géométriques : définition}

\begin{exo}[Coefficient multiplicateur]
\explique{Augmenter de $t\,\%$, c'est multiplier par $1+\dfrac{t}{100}$ ; diminuer de
$t\,\%$, c'est multiplier par $1-\dfrac{t}{100}$.}{rappel de seconde}
\explique{\etiq{a)} Augmenter de $4\,\%$ : $q=1+0{,}04=\res{1{,}04}$.}{$\dfrac{4}{100}=0{,}04$}
\explique{\etiq{b)} Diminuer de $4\,\%$ : $q=1-0{,}04=\res{0{,}96}$.}{}
\explique{\etiq{c)} Augmenter de $30\,\%$ : $q=1+0{,}3=\res{1{,}3}$.}{$\dfrac{30}{100}=0{,}3$}
\explique{\etiq{d)} Diminuer de $30\,\%$ : $q=1-0{,}3=\res{0{,}7}$.}{}
\explique{\etiq{e)} Augmenter de $60\,\%$ : $q=1+0{,}6=\res{1{,}6}$.}{}
\explique{\etiq{f)} Diminuer de $25\,\%$ : $q=1-0{,}25=\res{0{,}75}$.}{}
\explique{\etiq{g)} Augmenter de $2\,\%$ : $q=1+0{,}02=\res{1{,}02}$.}{$\dfrac{2}{100}=0{,}02$}
\explique{\etiq{h)} Diminuer de $45\,\%$ : $q=1-0{,}45=\res{0{,}55}$.}{}
\rem{Erreurs fréquentes : $4\,\%$ donne $0{,}04$, pas $0{,}4$ ; diminuer de $30\,\%$ donne
$0{,}7$, pas $0{,}3$ (multiplier par $0{,}3$, c'est ne garder que $30\,\%$).}
\end{exo}

\begin{exo}[Retrouver l'évolution]
\explique{Si $q>1$, c'est une hausse de $q-1$ ; si $q<1$, c'est une baisse de $1-q$. On
écrit ensuite ce nombre en pourcentage.}{$0{,}04$, c'est $4\,\%$}
\explique{\etiq{a)} $q=0{,}96<1$ : $1-0{,}96=0{,}04$, \res{diminuer de $4\,\%$}.}{}
\explique{\etiq{b)} $q=1{,}06>1$ : $1{,}06-1=0{,}06$, \res{augmenter de $6\,\%$}.}{}
\explique{\etiq{c)} $q=0{,}4<1$ : $1-0{,}4=0{,}6$, \res{diminuer de $60\,\%$}.}{on garde
$40\,\%$, on perd $60\,\%$}
\explique{\etiq{d)} $q=1{,}3>1$ : $1{,}3-1=0{,}3$, \res{augmenter de $30\,\%$}.}{}
\explique{\etiq{e)} $q=3>1$ : $3-1=2$, \res{augmenter de $200\,\%$}.}{tripler, c'est
ajouter deux fois la valeur de départ}
\explique{\etiq{f)} $q=0{,}98<1$ : $1-0{,}98=0{,}02$, \res{diminuer de $2\,\%$}.}{}
\explique{\etiq{g)} $q=0{,}85<1$ : $1-0{,}85=0{,}15$, \res{diminuer de $15\,\%$}.}{}
\explique{\etiq{h)} $q=1{,}075>1$ : $1{,}075-1=0{,}075$, \res{augmenter de
$7{,}5\,\%$}.}{$0{,}075\times 100=7{,}5$}
\rem{Erreur fréquente : lire $q=0{,}4$ comme « diminuer de $40\,\%$ ». Multiplier par
$0{,}4$, c'est garder $40\,\%$, donc perdre $60\,\%$.}
\end{exo}

\begin{exo}[Appliquer une évolution]
\explique{On multiplie par le coefficient multiplicateur. De tête, on peut aussi
calculer la variation, puis l'ajouter ou la retirer.}{les deux façons donnent le même
résultat}
\cas{a)}
\begin{calculs}
300\times 1{,}2 &= \res{360} &&\justif{$+20\,\%$ : $\times 1{,}2$ ; ou $300+60$, car
$20\,\%$ de $300$ font $60$}
\end{calculs}
\cas{b)}
\begin{calculs}
500\times 0{,}9 &= \res{450} &&\justif{$-10\,\%$ : $\times 0{,}9$ ; ou $500-50$}
\end{calculs}
\cas{c)}
\begin{calculs}
2\,000\times 1{,}03 &= \res{2\,060} &&\justif{$3\,\%$ de $2\,000$ font $60$ ; $2\,000+60$}
\end{calculs}
\cas{d)}
\begin{calculs}
60\times 0{,}75 &= \res{45} &&\justif{$25\,\%$, c'est un quart : $60-15$}
\end{calculs}
\end{exo}

\begin{exo}[Calculer les premiers termes]
\cas{a)}
\begin{calculs}
u_{1} &= 2\times 5 = \res{10} &&\justif{géométrique : on multiplie par la raison $2$}\\
u_{2} &= 2\times 10 = \res{20} &&\\
u_{3} &= 2\times 20 = \res{40} &&
\end{calculs}
\cas{b)}
\explique{Placer à $2\,\%$ par an, c'est augmenter de $2\,\%$ par an : on multiplie par
$1{,}02$.}{$1+\dfrac{2}{100}=1{,}02$}
\begin{calculs}
c_{0} &= \res{5\,000} &&\justif{le capital placé}\\
c_{1} &= 1{,}02\times 5\,000 = \res{5\,100} &&\justif{$2\,\%$ de $5\,000$ font $100$}\\
c_{2} &= 1{,}02\times 5\,100 = \res{5\,202} &&\justif{$2\,\%$ de $5\,100$ font $102$}
\end{calculs}
\explique{$c_{n+1}=\res{1{,}02\,c_{n}}$}{même coefficient chaque année}
\cas{c)}
\explique{$5\,\%$ de pertes, c'est multiplier par $0{,}95$.}{$1-\dfrac{5}{100}=0{,}95$}
\begin{calculs}
f_{0} &= \res{40\,000} &&\\
f_{1} &= 0{,}95\times 40\,000 = \res{38\,000} &&\justif{$5\,\%$ de $40\,000$ font $2\,000$}\\
f_{2} &= 0{,}95\times 38\,000 = \res{36\,100} &&\justif{$5\,\%$ de $38\,000$ font $1\,900$}
\end{calculs}
\explique{$f_{n+1}=\res{0{,}95\,f_{n}}$}{}
\cas{d)}
\explique{Baisser de $10\,\%$, c'est multiplier par $0{,}9$.}{$1-\dfrac{10}{100}=0{,}9$}
\begin{calculs}
e_{0} &= \res{800} &&\\
e_{1} &= 0{,}9\times 800 = \res{720} &&\justif{$10\,\%$ de $800$ font $80$}\\
e_{2} &= 0{,}9\times 720 = \res{648} &&\justif{$10\,\%$ de $720$ font $72$}
\end{calculs}
\explique{$e_{n+1}=\res{0{,}9\,e_{n}}$}{}
\cas{e)}
\explique{Augmenter de $50\,\%$, c'est multiplier par $1{,}5$.}{$1+\dfrac{50}{100}=1{,}5$}
\begin{calculs}
i_{0} &= \res{300} &&\\
i_{1} &= 1{,}5\times 300 = \res{450} &&\justif{$50\,\%$, c'est la moitié : $300+150$}\\
i_{2} &= 1{,}5\times 450 = \res{675} &&\justif{$450+225$}
\end{calculs}
\explique{$i_{n+1}=\res{1{,}5\,i_{n}}$}{}
\rem{Le pourcentage s'applique au terme précédent, pas au premier : $c_{2}=5\,202$, et non
$5\,200$.}
\end{exo}

\begin{exo}[Reconnaître une suite géométrique]
\explique{Pour une liste, on calcule les quotients « terme suivant divisé par terme
précédent » ; s'ils sont tous égaux, la suite est géométrique, de raison ce quotient. Pour
une situation, on cherche le même pourcentage à chaque étape.}{propriété 5}
\cas{a)}
\explique{$\dfrac{12}{3}=4$ ; \ $\dfrac{48}{12}=4$ ; \ $\dfrac{192}{48}=4$.}{quotients
successifs}
\explique{Géométrique, de raison \res{$q=4$}.}{les quotients sont égaux}
\cas{b)}
\explique{$\dfrac{32}{64}=0{,}5$ ; \ $\dfrac{16}{32}=0{,}5$ ; \ $\dfrac{8}{16}=0{,}5$.}{on
prend la moitié à chaque fois}
\explique{Géométrique, de raison \res{$q=0{,}5$}.}{diviser par $2$, c'est multiplier par
$0{,}5$}
\cas{c)}
\explique{$\dfrac{8}{4}=2$ mais $\dfrac{12}{8}=1{,}5$.}{les quotients changent}
\explique{\res{Pas géométrique}.}{elle est arithmétique, de raison $4$}
\cas{d)}
\explique{$\dfrac{550}{500}=1{,}1$. On vérifie : $550\times 1{,}1=605$ et
$605\times 1{,}1=665{,}5$.}{$\times 1{,}1$ : on ajoute $10\,\%$ ; $550+55$ et $605+60{,}5$}
\explique{Géométrique, de raison \res{$q=1{,}1$}.}{}
\cas{e)}
\explique{Augmenter de $4\,\%$, c'est multiplier par $1{,}04$ chaque année : suite
géométrique, de raison \res{$q=1{,}04$}.}{même pourcentage chaque année}
\cas{f)}
\explique{On ajoute $15$~€ chaque année : la même quantité, pas le même pourcentage.
\res{Pas géométrique}.}{elle est arithmétique, de raison $15$}
\end{exo}

\partie{Suites géométriques : expression en fonction de $n$}

\begin{exo}[Appliquer la formule]
\explique{On utilise $u_{n}=u_{0}\times q^{n}$ : il faut le premier terme $u_{0}$ et la
raison $q$.}{propriété 6}
\cas{a)}
\begin{calculs}
u_{n} &= u_{0}\times q^{n} &&\justif{premier terme $2$, raison $3$}\\
u_{n} &= \res{2\times 3^{n}} &&\\
u_{4} &= 2\times 3^{4} &&\justif{on remplace $n$ par $4$}\\
u_{4} &= 2\times 81 &&\justif{la puissance d'abord : $3\times 3\times 3\times 3=81$}\\
u_{4} &= \res{162} &&
\end{calculs}
\cas{b)}
\begin{calculs}
u_{n} &= \res{5\,000\times 0{,}1^{n}} &&\justif{premier terme $5\,000$, raison $0{,}1$}\\
u_{3} &= 5\,000\times 0{,}1^{3} &&\\
u_{3} &= 5\,000\times 0{,}001 &&\justif{$0{,}1\times 0{,}1\times 0{,}1=0{,}001$}\\
u_{3} &= \res{5} &&\justif{multiplier par $0{,}001$, c'est diviser par $1\,000$}
\end{calculs}
\cas{c)}
\explique{Le premier terme donné est $u_{1}$ : on recule d'un rang.}{la formule part de
$u_{0}$}
\begin{calculs}
u_{0} &= \dfrac{u_{1}}{q} &&\justif{on recule d'un rang : on divise par la raison}\\
u_{0} &= \dfrac{6}{2} = 3 &&\\
u_{n} &= \res{3\times 2^{n}} &&\justif{$u_{0}\times q^{n}$}\\
u_{5} &= 3\times 2^{5} &&\\
u_{5} &= 3\times 32 &&\justif{$2^{5}=2\times 2\times 2\times 2\times 2=32$}\\
u_{5} &= \res{96} &&
\end{calculs}
\explique{Contrôle : $u_{1}=3\times 2=6$.}{on retrouve la donnée}
\cas{d)}
\begin{calculs}
u_{n} &= \res{128\times 0{,}5^{n}} &&\justif{premier terme $128$, raison $0{,}5$}\\
u_{4} &= 128\times 0{,}5^{4} &&\\
u_{4} &= \dfrac{128}{2^{4}} &&\justif{multiplier $4$ fois par $0{,}5$, c'est diviser
$4$ fois par $2$}\\
u_{4} &= \dfrac{128}{16} &&\\
u_{4} &= \res{8} &&
\end{calculs}
\explique{Contrôle : $128$ ; $64$ ; $32$ ; $16$ ; $8$.}{on prend la moitié quatre fois}
\rem{Erreur fréquente : calculer $(2\times 3)^{4}=6^{4}$. La puissance ne porte que sur la
raison : $2\times 3^{4}=2\times 81$.}
\end{exo}

\begin{exo}[Placer ses économies]
\cas{a)}
\begin{calculs}
c_{1} &= 1{,}02\times 1\,500 &&\justif{$+2\,\%$ : on multiplie par $1{,}02$}\\
c_{1} &= \res{1\,530} &&\justif{$2\,\%$ de $1\,500$ font $30$}\\
c_{2} &= 1{,}02\times 1\,530 &&\\
c_{2} &= \res{1\,560{,}60} &&\justif{$2\,\%$ de $1\,530$ font $30{,}60$}
\end{calculs}
\explique{Au bout de deux ans, le capital vaut $1\,560{,}60$~€.}{}
\cas{b)}
\explique{Le capital augmente du même pourcentage chaque année : $(c_{n})$ est
\res{géométrique}, de premier terme $c_{0}=1\,500$ et de raison $q=1{,}02$.}{on multiplie
toujours par $1{,}02$}
\begin{calculs}
c_{n} &= c_{0}\times q^{n} &&\\
c_{n} &= \res{1\,500\times 1{,}02^{n}} &&
\end{calculs}
\cas{c)}
\begin{calculs}
c_{10} &= 1\,500\times 1{,}02^{10} &&\justif{on remplace $n$ par $10$}\\
c_{10} &\approx 1\,500\times 1{,}22 &&\justif{on utilise la valeur donnée}\\
c_{10} &\approx \res{1\,830} &&\justif{$1\,500+1\,500\times 0{,}22=1\,500+330$}
\end{calculs}
\explique{Au bout de $10$ ans, le capital vaut environ $1\,830$~€.}{« environ » :
$1{,}22$ est une valeur arrondie}
\end{exo}

\begin{exo}[Covoiturage]
\cas{a)}
\explique{Augmenter de $10\,\%$, c'est multiplier par $1{,}1$ chaque année : $(u_{n})$
est \res{géométrique}, de premier terme $u_{0}=80$ et de raison
\res{$q=1{,}1$}.}{$1+\dfrac{10}{100}=1{,}1$}
\cas{b)}
\explique{En C2 se trouve $u_{1}$ (2025), qui s'obtient en multipliant $u_{0}$, en B2 juste
à gauche, par $1{,}1$.}{les termes sont en ligne : on utilise la cellule de gauche}
\explique{En C2 : \res{\texttt{=B2*1,1}}}{recopiée vers la droite, elle devient
\texttt{=C2*1,1} en D2}
\cas{c)}
\begin{calculs}
u_{n} &= u_{0}\times q^{n} &&\\
u_{n} &= \res{80\times 1{,}1^{n}} &&
\end{calculs}
\cas{d)}
\begin{calculs}
u_{3} &= 80\times 1{,}1^{3} &&\justif{on remplace $n$ par $3$}\\
u_{3} &= 80\times 1{,}331 &&\justif{valeur donnée}\\
u_{3} &= 106{,}48 &&\justif{$8\times 1{,}331=10{,}648$, puis $\times 10$}
\end{calculs}
\explique{À l'unité près, $u_{3}\approx\res{106}$ : environ $106$ salariés en 2027.}{$n=3$ :
$2024+3=2027$}
\cas{e)}
\explique{Un quart des $400$ salariés : $\dfrac{400}{4}=100$ salariés.}{l'objectif}
\explique{En 2027, environ $106$ salariés covoiturent, et $106\ge 100$ : l'objectif
\res{sera atteint}.}{2027 correspond à $n=3$, calculé en d)}
\end{exo}

\begin{exo}[Dioxyde d'azote]
\cas{a)}
\begin{calculs}
u_{1} &= 0{,}95\times 60 &&\justif{baisser de $5\,\%$, c'est multiplier par $0{,}95$}\\
u_{1} &= \res{57} &&\justif{$5\,\%$ de $60$ font $3$ ; $60-3$}
\end{calculs}
\explique{Chaque année, on multiplie par $0{,}95$ : $(u_{n})$ est géométrique, de premier
terme $u_{0}=60$ et de raison $q=0{,}95$.}{même pourcentage chaque année}
\begin{calculs}
u_{n} &= u_{0}\times q^{n} &&\\
u_{n} &= \res{60\times 0{,}95^{n}} &&
\end{calculs}
\cas{b)}
\begin{calculs}
u_{6} &= 60\times 0{,}95^{6} &&\justif{2026 correspond à $n=6$}\\
u_{6} &\approx 60\times 0{,}74 &&\justif{valeur donnée}\\
u_{6} &\approx \res{44{,}4} &&\justif{$6\times 0{,}74=4{,}44$, puis $\times 10$}
\end{calculs}
\explique{En 2026, la concentration sera d'environ $44{,}4$~µg/m$^{3}$.}{}
\end{exo}

\begin{exo}[Lire la formule]
\explique{Une suite géométrique s'écrit $u_{n}=u_{0}\times q^{n}$ : le nombre devant est le
premier terme, le nombre élevé à la puissance $n$ est la raison.}{on compare avec la
formule}
\cas{a)}
\explique{$u_{0}=\res{4}$ et $q=\res{5}$.}{contrôle : $u_{0}=4\times 5^{0}=4$}
\cas{b)}
\explique{$v_{0}=\res{2\,000}$ et $q=\res{0{,}85}$.}{}
\cas{c)}
\explique{$w_{n}=3^{n}$ s'écrit $w_{n}=1\times 3^{n}$ : $w_{0}=\res{1}$ et
$q=\res{3}$.}{il n'y a pas de nombre devant : c'est $1$, et non $0$}
\end{exo}

\partie{Suites géométriques : variations et représentation graphique}

\begin{exo}[Sens de variation]
\explique{Il suffit de la raison (le premier terme est positif) : si $q>1$, la suite est
croissante ; si $0<q<1$, elle est décroissante.}{propriété 7}
\cas{a)}
\explique{$q=\res{1{,}3}$ et $1{,}3>1$ : $(u_{n})$ est \res{croissante}.}{la raison est le
nombre élevé à la puissance $n$}
\cas{b)}
\explique{$q=\res{0{,}6}$ et $0{,}6<1$ : $(v_{n})$ est \res{décroissante}.}{}
\cas{c)}
\explique{$q=\res{1{,}02}$ et $1{,}02>1$ : $(w_{n})$ est \res{croissante}.}{on multiplie
par $1{,}02$ à chaque étape}
\cas{d)}
\explique{$q=\res{0{,}25}$ et $0{,}25<1$ : $(t_{n})$ est \res{décroissante}.}{}
\cas{e)}
\explique{Baisser de $5\,\%$, c'est multiplier par $0{,}95$ : $q=\res{0{,}95}$, et
$0{,}95<1$ : la suite est \res{décroissante}.}{$1-\dfrac{5}{100}=0{,}95$}
\rem{Erreur fréquente : regarder le premier terme ($8$, $500$\ldots) au lieu de la raison.}
\end{exo}

\begin{exo}[Linéaire ou exponentielle ?]
\cas{a)}
\explique{Croix droites, $(a_{n})$ : $8$ ; $6{,}5$ ; $5$ ; $3{,}5$ ; $2$.}{on lit les
ordonnées}
\explique{Différences : $6{,}5-8=-1{,}5$, et ainsi de suite, toujours $-1{,}5$.}{on retire
$1{,}5$ à chaque rang}
\explique{$(a_{n})$ est \res{arithmétique}, de raison \res{$r=-1{,}5$}.}{}
\explique{Croix obliques, $(b_{n})$ : $8$ ; $4$ ; $2$ ; $1$ ; $0{,}5$.}{}
\explique{Quotients : $\dfrac{4}{8}=0{,}5$, et ainsi de suite, toujours $0{,}5$.}{on prend
la moitié à chaque rang}
\explique{$(b_{n})$ est \res{géométrique}, de raison \res{$q=0{,}5$}.}{}
\cas{b)}
\begin{calculs}
a_{5} &= a_{4}+r = 2-1{,}5 &&\justif{on avance d'un rang : on ajoute $-1{,}5$}\\
a_{5} &= \res{0{,}5} &&\\
b_{5} &= b_{4}\times q = 0{,}5\times 0{,}5 &&\justif{on multiplie par $0{,}5$}\\
b_{5} &= \res{0{,}25} &&
\end{calculs}
\cas{c)}
\explique{$(a_{n})$ : points alignés, décroissance \res{linéaire}.}{suite arithmétique}
\explique{$(b_{n})$ : points non alignés, qui descendent de moins en moins vite :
décroissance \res{exponentielle}.}{suite géométrique de raison $0<q<1$}
\end{exo}

\begin{exo}[Vrai ou faux ?]
\cas{a)}
\explique{\res{Faux} : $q=1{,}05>1$, la suite est croissante.}{propriété 7}
\cas{b)}
\explique{\res{Vrai} : baisser de $4\,\%$, c'est multiplier par $q=0{,}96$, et
$0{,}96<1$ : la suite est décroissante.}{$1-\dfrac{4}{100}=0{,}96$}
\cas{c)}
\explique{\res{Vrai} : les points d'une suite arithmétique sont toujours alignés, quelle
que soit la raison.}{propriété 4}
\cas{d)}
\explique{\res{Faux} : $5$ est le premier terme ; la raison est $0{,}8$, le nombre élevé à
la puissance $n$. Comme $0{,}8<1$, la suite est décroissante.}{on compare la raison à $1$,
pas le premier terme}
\end{exo}

\partie{Suites géométriques : problème de seuil}

\begin{exo}[Dioxyde d'azote (suite)]
\cas{a)}
\begin{calculs}
60\times 0{,}95^{n} &\le 45 &&\justif{on traduit $u_{n}\le 45$}\\
0{,}95^{n} &\le \dfrac{45}{60} &&\justif{on divise par $60$, positif : le sens ne change
pas}\\
0{,}95^{n} &\le 0{,}75 &&\justif{$\dfrac{45}{60}=\dfrac{3\times\barre{15}}{4\times\barre{15}}$,
soit $\dfrac{3}{4}$}
\end{calculs}
\cas{b)}
\explique{$0{,}95^{5}\approx 0{,}77$ et $0{,}77>0{,}75$ : pour $n=5$, l'objectif n'est pas
atteint.}{on compare chaque valeur donnée à $0{,}75$}
\explique{$0{,}95^{6}\approx 0{,}74$ et $0{,}74\le 0{,}75$ : le premier rang qui convient
est $n=6$.}{}
\explique{L'objectif sera atteint \res{en 2026}.}{$2020+6=2026$ ; on retrouve
$u_{6}\approx 44{,}4$ (exercice 40)}
\rem{La valeur $0{,}95^{7}$ ne sert pas : on cherche le premier rang qui convient. Comme la
suite est décroissante, l'objectif reste atteint ensuite.}
\end{exo}

\begin{exo}[Lettre d'information]
\cas{a)}
\explique{Augmenter de $6\,\%$, c'est multiplier par $1{,}06$ chaque mois : $(u_{n})$ est
géométrique, de premier terme $u_{0}=2\,500$ et de raison $q=1{,}06$.}{$1+\dfrac{6}{100}=1{,}06$}
\begin{calculs}
u_{n} &= u_{0}\times q^{n} &&\\
u_{n} &= \res{2\,500\times 1{,}06^{n}} &&
\end{calculs}
\cas{b)}
\begin{calculs}
2\,500\times 1{,}06^{n} &\ge 5\,000 &&\justif{on traduit $u_{n}\ge 5\,000$}\\
1{,}06^{n} &\ge \dfrac{5\,000}{2\,500} &&\justif{on divise par $2\,500$, positif}\\
1{,}06^{n} &\ge 2 &&
\end{calculs}
\cas{c)}
\explique{« Doubler », c'est passer de $2\,500$ à $5\,000$ abonnés : on cherche le premier
$n$ tel que $1{,}06^{n}\ge 2$.}{question b)}
\explique{$1{,}06^{11}\approx 1{,}90$ et $1{,}90<2$ : pas encore.}{}
\explique{$1{,}06^{12}\approx 2{,}01$ et $2{,}01\ge 2$ : le premier rang est $n=12$.}{}
\explique{Le nombre d'abonnés aura doublé \res{au bout de $12$ mois}, en janvier de
l'année suivante.}{on répond dans le contexte}
\end{exo}

\begin{exo}[La valeur d'un scooter]
\cas{a)}
\explique{Perdre $20\,\%$, c'est multiplier par $0{,}8$ chaque année : $(v_{n})$ est
géométrique, de premier terme $v_{0}=3\,000$ et de raison $q=0{,}8$.}{$1-\dfrac{20}{100}=0{,}8$}
\begin{calculs}
v_{n} &= v_{0}\times q^{n} &&\\
v_{n} &= \res{3\,000\times 0{,}8^{n}} &&
\end{calculs}
\cas{b)}
\begin{calculs}
3\,000\times 0{,}8^{n} &< 1\,200 &&\justif{on traduit $v_{n}<1\,200$}\\
0{,}8^{n} &< \dfrac{1\,200}{3\,000} &&\justif{on divise par $3\,000$, positif}\\
0{,}8^{n} &< 0{,}4 &&\justif{$\dfrac{12}{30}=\dfrac{4}{10}$}
\end{calculs}
\explique{$0{,}8^{4}\approx 0{,}41$ et $0{,}41>0{,}4$ : pas encore.}{on compare chaque
valeur donnée à $0{,}4$}
\explique{$0{,}8^{5}\approx 0{,}33$ et $0{,}33<0{,}4$ : le premier rang est $n=5$.}{}
\explique{Le scooter vaudra moins de $1\,200$~€ \res{au bout de $5$ ans}.}{}
\rem{Piège : $0{,}41$ est très proche de $0{,}4$, mais plus grand ; $n=4$ ne convient
pas.}
\end{exo}

\partie{Comparer deux évolutions}

\begin{exo}[Arithmétique, géométrique ou aucune des deux ?]
\explique{Formule $u_{0}+r\times n$ ou relation « on ajoute $r$ » : arithmétique. Formule
$u_{0}\times q^{n}$ ou relation « on multiplie par $q$ » : géométrique.}{tableau-bilan du
cours}
\cas{a)}
\explique{$u_{n}=5n$ s'écrit $0+5n$ : \res{arithmétique}, de raison \res{$r=5$}.}{on
ajoute $5$ à chaque rang}
\cas{b)}
\explique{$u_{0}=-1$ ; \ $u_{1}=0$ ; \ $u_{2}=3$.}{on calcule les premiers termes}
\explique{Différences : $1$, puis $3$ : pas arithmétique. On ne passe pas de $0$ à $3$ en
multipliant : pas géométrique.}{le carré de $n$ : ni « $+r$ » ni « $\times q$ »}
\explique{\res{Ni l'une ni l'autre}.}{}
\cas{c)}
\explique{$u_{n}=4\times 3^{n}$ : \res{géométrique}, de premier terme $4$ et de raison
\res{$q=3$}.}{forme $u_{0}\times q^{n}$}
\cas{d)}
\explique{$u_{n+1}=u_{n}-2$ : on ajoute $-2$ à chaque étape, \res{arithmétique}, de raison
\res{$r=-2$}.}{forme $u_{n}+r$}
\cas{e)}
\begin{calculs}
u_{n+1} &= u_{n}+0{,}2\,u_{n} &&\justif{on recopie la relation}\\
u_{n+1} &= 1\times u_{n}+0{,}2\times u_{n} &&\justif{$u_{n}$, c'est $1\times u_{n}$}\\
u_{n+1} &= 1{,}2\,u_{n} &&\justif{$1+0{,}2=1{,}2$}
\end{calculs}
\explique{\res{Géométrique}, de raison \res{$q=1{,}2$}.}{on multiplie par $1{,}2$ :
hausse de $20\,\%$}
\cas{f)}
\explique{$u_{n+1}=2u_{n}$ : \res{géométrique}, de raison \res{$q=2$}.}{forme $q\times
u_{n}$}
\rem{e) est un piège : on ajoute $0{,}2\,u_{n}$, qui change à chaque rang, et non un nombre
fixe. Ajouter $0{,}2\,u_{n}$, c'est augmenter de $20\,\%$.}
\end{exo}

\begin{exo}[Deux placements]
\cas{a)}
\begin{calculs}
a_{1} &= 15\,000+400 = \res{15\,400} &&\justif{placement A : on ajoute $400$~€}\\
b_{1} &= 1{,}03\times 15\,000 = \res{15\,450} &&\justif{placement B : $3\,\%$ de
$15\,000$ font $450$}
\end{calculs}
\explique{Au bout d'un an, le placement A donne $15\,400$~€, le placement B
$15\,450$~€.}{}
\cas{b)}
\explique{$(a_{n})$ : on ajoute $400$ chaque année, suite \res{arithmétique}, de raison
\res{$r=400$}.}{même quantité ajoutée}
\explique{$(b_{n})$ : on multiplie par $1{,}03$ chaque année, suite \res{géométrique}, de
raison \res{$q=1{,}03$}.}{même pourcentage : $1+\dfrac{3}{100}$}
\cas{c)}
\begin{calculs}
a_{n} &= \res{15\,000+400n} &&\justif{$a_{0}+n\times r$}\\
b_{n} &= \res{15\,000\times 1{,}03^{n}} &&\justif{$b_{0}\times q^{n}$}
\end{calculs}
\cas{d)}
\begin{calculs}
15\,000+400n &\ge 18\,000 &&\justif{le projet coûte $18\,000$~€}\\
400n &\ge 3\,000 &&\justif{on retire $15\,000$}\\
n &\ge 7{,}5 &&\justif{on divise par $400$ : $\dfrac{3\,000}{400}=\dfrac{30}{4}$}
\end{calculs}
\explique{Le premier entier qui convient est $n=8$ : avec le placement A, le projet sera
financé \res{en 2033}.}{$2025+8=2033$ ; contrôle : $a_{7}=17\,800$ et $a_{8}=18\,200$}
\cas{e)}
\explique{$b_{6}\approx 17\,911$ et $17\,911<18\,000$ : pas encore.}{on compare les valeurs
données à $18\,000$}
\explique{$b_{7}\approx 18\,448$ et $18\,448\ge 18\,000$ : le premier rang est
$n=7$.}{}
\explique{Avec le placement B, le projet sera financé \res{en 2032}, un an plus tôt
qu'avec A : il vaut mieux choisir \res{le placement B}.}{$2025+7=2032$}
\end{exo}

\begin{exo}[Réduire ses déchets]
\cas{a)}
\explique{$(a_{n})$ : $15$~t de moins chaque année, suite \res{arithmétique}, de raison
\res{$r=-15$}.}{même quantité retirée}
\explique{$(b_{n})$ : $7\,\%$ de moins chaque année, on multiplie par $0{,}93$, suite
\res{géométrique}, de raison \res{$q=0{,}93$}.}{$1-\dfrac{7}{100}=0{,}93$}
\cas{b)}
\begin{calculs}
a_{n} &= \res{360-15n} &&\justif{premier terme $360$, raison $-15$}\\
b_{n} &= \res{500\times 0{,}93^{n}} &&\justif{premier terme $500$, raison $0{,}93$}
\end{calculs}
\cas{c)}
\explique{Objectif de A : la moitié de $360$, soit $180$~t au plus.}{« diviser par deux »}
\begin{calculs}
360-15n &\le 180 &&\justif{on traduit la question}\\
-15n &\le -180 &&\justif{on retire $360$}\\
n &\ge 12 &&\justif{on divise par $-15$, négatif : le sens change}
\end{calculs}
\explique{L'entreprise A atteindra son objectif \res{en 2032}.}{$2020+12=2032$ ;
contrôle : $a_{12}=360-180=180$}
\cas{d)}
\explique{Objectif de B : la moitié de $500$, soit $250$~t au plus.}{}
\begin{calculs}
500\times 0{,}93^{n} &\le 250 &&\justif{on traduit la question}\\
0{,}93^{n} &\le 0{,}5 &&\justif{on divise par $500$ : $\dfrac{250}{500}=0{,}5$}
\end{calculs}
\explique{$0{,}93^{9}\approx 0{,}52>0{,}5$ et $0{,}93^{10}\approx 0{,}48\le 0{,}5$ : le
premier rang est $n=10$.}{on compare les valeurs données à $0{,}5$}
\explique{L'entreprise B atteindra son objectif \res{en 2030}.}{$2020+10=2030$}
\cas{e)}
\explique{\res{L'entreprise B} y parvient la première, en 2030, deux ans avant
l'entreprise A.}{$2030$ contre $2032$}
\end{exo}

\partiesansnum{Exercices type épreuve anticipée}

\begin{exo}[Automatismes (QCM)]
\cas{1.}
\explique{Réponse \res{c)} : $1-\dfrac{30}{100}=0{,}7$.}{multiplier par $0{,}3$, ce serait
ne garder que $30\,\%$}
\cas{2.}
\explique{Réponse \res{b)} : $\dfrac{8}{4}=2$ ; \ $\dfrac{16}{8}=2$ ; \ $\dfrac{32}{16}=2$,
on multiplie toujours par $2$.}{les différences $4$ ; $8$ ; $16$ changent : pas
arithmétique}
\cas{3.}
\explique{Réponse \res{b)} : $u_{1}=10$ ; \ $u_{2}=20$ ; \ $u_{3}=40$.}{on double trois
fois}
\cas{4.}
\explique{Réponse \res{b)} : $u_{10}=60-3\times 10=60-30$, soit $30$.}{on remplace $n$
par $10$}
\cas{5.}
\explique{Réponse \res{a)} : la raison $1{,}02$ est plus grande que $1$, et le premier
terme est positif.}{propriété 7}
\cas{6.}
\explique{Réponse \res{a)} : $u_{1}=400\times 0{,}5$, soit $200$.}{$0{,}5^{1}=0{,}5$ :
la moitié de $400$}
\end{exo}

\begin{exo}[Location de vélos]
\cas{1.}
\explique{$30$ vélos de plus chaque mois : $(u_{n})$ est arithmétique, de premier terme
$u_{0}=40$ et de raison $r=30$.}{même quantité ajoutée chaque mois}
\begin{calculs}
u_{n} &= u_{0}+n\times r &&\\
u_{n} &= \res{40+30n} &&\\
u_{10} &= 40+30\times 10 &&\justif{au mois $10$ : $n=10$}\\
u_{10} &= 40+300 &&\\
u_{10} &= 340 &&
\end{calculs}
\explique{$340>300$ : au mois $10$, l'entreprise possède \res{plus de $300$ vélos}.}{}
\cas{2.}
\begin{calculs}
40+30n &= 640 &&\justif{on traduit $u_{n}=640$}\\
30n &= 600 &&\justif{on retire $40$}\\
n &= 20 &&\justif{on divise par $30$}
\end{calculs}
\explique{Elle aura $640$ vélos \res{au bout de $20$ mois}.}{contrôle :
$u_{20}=40+600=640$}
\cas{3.}
\explique{Augmenter de $20\,\%$, c'est multiplier par $1{,}2$ : d'un mois au suivant,
$v_{n+1}=1{,}2\,v_{n}$.}{$1+\dfrac{20}{100}=1{,}2$}
\explique{On multiplie toujours par le même nombre : $(v_{n})$ est \res{géométrique}, de
premier terme $v_{0}=25$ et de raison $q=1{,}2$.}{}
\cas{4.}
\begin{calculs}
v_{n} &= v_{0}\times q^{n} &&\\
v_{n} &= \res{25\times 1{,}2^{n}} &&
\end{calculs}
\cas{5.}
\begin{calculs}
v_{5} &= 25\times 1{,}2^{5} &&\justif{on remplace $n$ par $5$}\\
v_{5} &\approx 25\times 2{,}49 &&\justif{valeur donnée}\\
v_{5} &\approx 62{,}25 &&\justif{$25\times 2=50$ et $25\times 0{,}49=12{,}25$}
\end{calculs}
\explique{Donc $v_{5}\approx 62$ : au mois $5$, l'entreprise compte \res{environ $62$
clients}.}{un nombre de clients est entier}
\cas{6.}
\begin{calculs}
25\times 1{,}2^{n} &> 150 &&\justif{on traduit $v_{n}>150$}\\
1{,}2^{n} &> 6 &&\justif{on divise par $25$ : $\dfrac{150}{25}=6$}
\end{calculs}
\explique{$1{,}2^{9}\approx 5{,}16$ et $5{,}16<6$ : pas encore.}{on compare les valeurs
données à $6$}
\explique{$1{,}2^{10}\approx 6{,}19$ et $6{,}19>6$ : le premier rang est $n=10$.}{}
\explique{L'entreprise aura plus de $150$ clients \res{à partir du mois $10$}.}{}
\end{exo}

\begin{exo}[Deux clubs sportifs]
\cas{1.}
\begin{calculs}
T_{1} &= 700-15 = \res{685} &&\justif{on perd $15$ adhérents en un an}
\end{calculs}
\explique{$(T_{n})$ est arithmétique, de premier terme $T_{0}=700$ et de raison
$r=-15$.}{même nombre perdu chaque année}
\begin{calculs}
T_{n} &= T_{0}+n\times r &&\\
T_{n} &= \res{700-15n} &&
\end{calculs}
\cas{2.}
\begin{calculs}
T_{10} &= 700-15\times 10 &&\justif{2035 correspond à $n=10$}\\
T_{10} &= 700-150 &&\\
T_{10} &= 550 &&
\end{calculs}
\explique{Entre 2025 et 2035, le club perd $700-550=150$ adhérents.}{}
\begin{calculs}
0{,}2\times 700 &= 140 &&\justif{$20\,\%$ de $700$}
\end{calculs}
\explique{$150>140$ : le club de tennis aura bien perdu \res{plus de $20\,\%$} de ses
adhérents.}{on compare la perte à $20\,\%$ de l'effectif de départ}
\cas{3.}
\explique{2028 correspond à $n=3$. Le point d'abscisse $3$ est entre les lignes $200$ et
$300$, un peu au-dessus du milieu.}{une ligne horizontale tous les $100$ adhérents}
\explique{En 2028, le club d'escalade compte \res{environ $260$ adhérents}.}{une lecture
entre $250$ et $270$ convient}
\cas{4.}
\explique{$E_{n+1}=1{,}2\,E_{n}$ : on multiplie toujours par $1{,}2$, donc $(E_{n})$ est
\res{géométrique}, de premier terme \res{$E_{0}=150$} et de raison
\res{$q=1{,}2$}.}{$E_{0}$ : l'effectif de 2025}
\cas{5.}
\explique{$E_{n}=150\times 1{,}2^{n}$. On compare les deux clubs aux rangs $7$ et $8$, ceux
des valeurs données.}{$E_{0}\times q^{n}$}
\begin{calculs}
E_{7} &\approx 150\times 3{,}58 = 537 &&\justif{$100\times 3{,}58+50\times 3{,}58$, soit
$358+179$}\\
T_{7} &= 700-15\times 7 = 595 &&\justif{$700-105$}
\end{calculs}
\explique{$537<595$ : en 2032, le club de tennis a encore plus d'adhérents.}{}
\begin{calculs}
E_{8} &\approx 150\times 4{,}30 = 645 &&\justif{$430+215$}\\
T_{8} &= 700-15\times 8 = 580 &&\justif{$700-120$}
\end{calculs}
\explique{$645>580$ : en 2033, le club d'escalade passe devant.}{}
\explique{Comme $(E_{n})$ augmente et $(T_{n})$ diminue, l'escalade reste devant ensuite :
le club d'escalade aura plus d'adhérents que le club de tennis \res{à partir de
2033}.}{$q=1{,}2>1$ : croissante ; $r=-15<0$ : décroissante}
\end{exo}

\end{document}
